% ==============================================================================
% SECCIÓN 03.04: DISTRIBUCIONES DE BERNOULLI Y BINOMIAL
% ==============================================================================

\subsection*{Enunciados de los problemas --- Sección 03.04}

\subsubsection*{Nivel Fundamental}

\begin{problema}
	\label{prob:3.4.1}
	Considere un experimento aleatorio elemental consistente en la inspección de un único componente electrónico en una línea de montaje, donde la probabilidad de que el componente sea apto para su uso es $p = 0.85$. Sea $X$ la variable aleatoria que toma el valor $1$ si el componente es apto (éxito) y $0$ si resulta defectuoso (fracaso).
	\begin{enumerate}
		\item Especifique formalmente la función de masa de probabilidad $f_X(x)$ utilizando la forma compacta del ensayo de Bernoulli: $f_X(x) = p^x (1-p)^{1-x}$ para $x \in \set{0, 1}$.
		\item Demuestre directamente a partir de la definición que $\mathbb{E}[X] = p$ y $\text{Var}(X) = p(1-p)$.
	\end{enumerate}
\end{problema}

\begin{problema}
	\label{prob:3.4.2}
	En un proceso automático de control de calidad, se inspeccionan $n = 8$ sensores de temperatura fabricados en condiciones idénticas e independientes. La probabilidad de que un sensor individual cumpla con los estándares de precisión es $p = 0.90$. Sea $X$ el número total de sensores que cumplen con los estándares entre los $8$ inspeccionados.
	\begin{enumerate}
		\item Identifique la distribución exacta de $X$ junto con sus parámetros paramétricos y calcule la probabilidad de que exactamente $7$ sensores cumplan con los estándares.
		\item Calcule la probabilidad de que a lo sumo $1$ sensor resulte no apto en la muestra inspeccionada.
	\end{enumerate}
\end{problema}

\begin{problema}
	\label{prob:3.4.3}
	Sea un conjunto de ensayos independientes de Bernoulli $X_1, X_2, \dots, X_n$, donde cada variable aleatoria cumple que $P(X_i = 1) = p$ y $P(X_i = 0) = q = 1-p$ para todo $i \in \set{1, 2, \dots, n}$.
	\begin{enumerate}
		\item Defina la variable sumatoria $Y = \sum_{i=1}^n X_i$. Explique argumentativamente por qué el número de secuencias en el espacio muestral $\Omega$ que producen exactamente $k$ éxitos ($Y = k$) está dado por el coeficiente combinatorio $\comb{n}{k}$.
		\item Utilice este razonamiento para deducir de primeros principios la función de masa de probabilidad $P(Y = k) = \comb{n}{k} p^k q^{n-k}$ y verifique explícitamente su suma sobre el soporte $\mathcal{S}_Y = \set{0, 1, \dots, n}$.
	\end{enumerate}
\end{problema}

\subsubsection*{Nivel Operativo}

\begin{problema}
	\label{prob:3.4.4}
	Una firma de ciberseguridad monitorea una red de servidores donde la probabilidad de que un intento de acceso externo sea clasificado como malicioso es $p = 0.15$. En un intervalo de tiempo dado se registran $n = 12$ intentos de acceso independientes. Sea $X \sim \text{Bin}(12, 0.15)$ el número de accesos maliciosos detectados.
	\begin{enumerate}
		\item Calcule la probabilidad de que se detecten a lo sumo $2$ accesos maliciosos ($P(X \le 2)$).
		\item Calcule la probabilidad condicional de que se detecten al menos $4$ accesos maliciosos dado que se detectó al menos uno ($P(X \ge 4 \mid X \ge 1)$).
	\end{enumerate}
\end{problema}

\begin{problema}
	\label{prob:3.4.5}
	Un equipo de investigación en bioinformática desarrolla un algoritmo para la detección de una secuencia genética rara en muestras clínicas. La probabilidad de que una muestra positiva sea correctamente identificada en un ensayo individual es $p = 0.08$. Si los ensayos se realizan de manera totalmente independiente sobre muestras distintas:
	\begin{enumerate}
		\item Plantee la inecuación logarítmica que permite determinar el número mínimo de muestras independientes $n$ que deben procesarse para que la probabilidad de detectar la secuencia al menos una vez sea de mínimo $99\%$ ($P(X \ge 1) \ge 0.99$).
		\item Resuelva numéricamente para obtener el valor entero exacto de $n^*$.
	\end{enumerate}
\end{problema}

\begin{problema}
	\label{prob:3.4.6}
	Un fabricante de componentes aeroespaciales suministra lotes de $n = 20$ válvulas hidráulicas. La probabilidad de que una válvula defectuosa escape al control interno es $p = 0.05$. El contrato contractual establece un costo operativo de garantía modelado por la función de pérdida no lineal:
	\begin{align*}
		C(X) = 500 X^2 + 150 X + 1000 \quad \text{(en dólares)},
	\end{align*}
	donde $X$ representa el número de válvulas defectuosas en el lote entregado.
	\begin{enumerate}
		\item Determine la esperanza matemática del costo de garantía $\mathbb{E}[C(X)]$ utilizando las propiedades de los momentos binomiales.
		\item Calcule la desviación estándar del número de defectos $\sigma_X$ y discuta el impacto del término cuadrático $X^2$ en el costo esperado global.
	\end{enumerate}
\end{problema}

\subsubsection*{Nivel Analítico}

\begin{problema}
	\label{prob:3.4.7}
	Sea $X \sim \text{Bin}(n, p)$ con $q = 1-p$.
	\begin{enumerate}
		\item Demuestre rigurosamente, utilizando el Teorema del Binomio de Newton, la identidad de normalización absoluta:
		\begin{align*}
			\sum_{k=0}^n \comb{n}{k} p^k q^{n-k} = 1.
		\end{align*}
		\item Aplicando el operador diferencial $\frac{\partial}{\partial p}$ a la identidad analítica $(p+q)^n = \sum_{k=0}^n \comb{n}{k} p^k q^{n-k}$ antes de sustituir $q=1-p$, demuestre analíticamente que $\mathbb{E}[X] = np$.
	\end{enumerate}
\end{problema}

\begin{problema}
	\label{prob:3.4.8}
	Para una variable aleatoria discreta $X \sim \text{Bin}(n, p)$, el segundo momento factorial se define como el valor esperado del producto descendente $\mathbb{E}[X(X-1)]$.
	\begin{enumerate}
		\item Deduzca analíticamente la fórmula cerrada de $\mathbb{E}[X(X-1)]$ manipulando los coeficientes combinatorios de la suma $\sum_{k=2}^n k(k-1)\comb{n}{k} p^k q^{n-k}$.
		\item A partir de la identidad algebraica $\text{Var}(X) = \mathbb{E}[X(X-1)] + \mathbb{E}[X] - (\mathbb{E}[X])^2$, demuestre de forma exacta que $\text{Var}(X) = np(1-p)$.
	\end{enumerate}
\end{problema}

\subsubsection*{Nivel Desafiante}

\begin{problema}
	\label{prob:3.4.9}
	Sea $X \sim \text{Bin}(n, p)$ una variable binomial con probabilidad de éxito $0 < p < 1$. Se define la moda de la distribución como el valor o los valores $k^* \in \set{0, 1, \dots, n}$ que maximizan la probabilidad $P(X = k^*)$.
	\begin{enumerate}
		\item Analice el cociente de probabilidades adyacentes $R(k) = \dfrac{P(X = k)}{P(X = k-1)}$ para todo $k \in \set{1, 2, \dots, n}$.
		\item Demuestre analíticamente que la moda de la distribución cumple que $k^* = \lfloor (n+1)p \rfloor$ cuando $(n+1)p$ no es un número entero.
		\item Determine exactamente qué ocurre con la distribución cuando $(n+1)p = m \in \mathbb{Z}^+$ (caso bimodal).
	\end{enumerate}
\end{problema}

\begin{problema}
	\label{prob:3.4.10}
	Considere dos variables aleatorias binomiales mutuamente independientes: $X \sim \text{Bin}(n_1, p)$ y $Y \sim \text{Bin}(n_2, p)$, ambas compartiendo exactamente la misma probabilidad de éxito $p$ en cada ensayo.
	\begin{enumerate}
		\item Demuestre combinatoriamente, utilizando la convolución discreta y la Identidad de Vandermonde, que la suma $Z = X + Y$ sigue una distribución exacta $\text{Bin}(n_1 + n_2, p)$.
		\item Compruebe que la aditividad de la esperanza y la varianza se conserva de manera exacta bajo esta suma de distribuciones independientes.
	\end{enumerate}
\end{problema}

\subsection*{Sugerencias de los problemas --- Sección 03.04}

\begin{sugerencia}
	\label{sug:3.4.1}
	Para el Problema \ref{prob:3.4.1}, evalúe la expresión compacta $p^x(1-p)^{1-x}$ en $x=0$ y $x=1$ para verificar que coincide con la definición de probabilidad de éxito y fracaso. Para los momentos, aplique la suma directa $\mathbb{E}[X^k] = 0^k \cdot (1-p) + 1^k \cdot p$.
\end{sugerencia}

\begin{sugerencia}
	\label{sug:3.4.2}
	Para el Problema \ref{prob:3.4.2}, recuerde que la condición "a lo sumo $1$ sensor no apto" equivale exactamente a que al menos $7$ sensores cumplan con los estándares ($X \ge 7$), lo cual implica evaluar $P(X=7) + P(X=8)$.
\end{sugerencia}

\begin{sugerencia}
	\label{sug:3.4.3}
	Para el Problema \ref{prob:3.4.3}, cada secuencia particular de $k$ éxitos y $n-k$ fracasos tiene una probabilidad individual de $p^k q^{n-k}$ por independencia. El número total de configuraciones distintas para ubicar los $k$ éxitos en las $n$ posiciones es $\comb{n}{k}$.
\end{sugerencia}

\begin{sugerencia}
	\label{sug:3.4.4}
	Para el Problema \ref{prob:3.4.4}, en la parte condicional utilice la definición de probabilidad condicional: $P(X \ge 4 \mid X \ge 1) = \dfrac{P(X \ge 4 \cap X \ge 1)}{P(X \ge 1)} = \dfrac{1 - P(X \le 3)}{1 - P(X = 0)}$.
\end{sugerencia}

\begin{sugerencia}
	\label{sug:3.4.5}
	Para el Problema \ref{prob:3.4.5}, el evento complementario a "detectar al menos una vez" es "cero detecciones en $n$ ensayos" ($X=0$), cuya probabilidad es exactamente $(1-p)^n$. Exprese la condición como $1 - (1-p)^n \ge 0.99$.
\end{sugerencia}

\begin{sugerencia}
	\label{sug:3.4.6}
	Para el Problema \ref{prob:3.4.6}, por linealidad del operador esperanza, $\mathbb{E}[C(X)] = 500 \mathbb{E}[X^2] + 150 \mathbb{E}[X] + 1000$. Recuerde que el segundo momento bruto se obtiene como $\mathbb{E}[X^2] = \text{Var}(X) + (\mathbb{E}[X])^2 = np(1-p) + n^2 p^2$.
\end{sugerencia}

\begin{sugerencia}
	\label{sug:3.4.7}
	Para el Problema \ref{prob:3.4.7}, al derivar $(p+q)^n$ respecto a $p$ obtenemos $n(p+q)^{n-1}$. Al derivar el lado derecho de la suma obtenemos $\sum k \comb{n}{k} p^{k-1} q^{n-k}$. Multiplique por $p$ a ambos lados y luego sustituya $q = 1-p$.
\end{sugerencia}

\begin{sugerencia}
	\label{sug:3.4.8}
	Para el Problema \ref{prob:3.4.8}, note que para $k \ge 2$, $k(k-1)\comb{n}{k} = n(n-1)\comb{n-2}{k-2}$. Factorice $n(n-1)p^2$ fuera de la suma y reconozca la normalización de una binomial con parámetros $n-2$ y $p$.
\end{sugerencia}

\begin{sugerencia}
	\label{sug:3.4.9}
	Para el Problema \ref{prob:3.4.9}, plantee la condición de crecimiento $R(k) = \dfrac{\comb{n}{k} p^k q^{n-k}}{\comb{n}{k-1} p^{k-1} q^{n-k+1}} = \dfrac{(n-k+1)p}{kq} \ge 1$. Despeje la variable $k$ en función de $(n+1)p$.
\end{sugerencia}

\begin{sugerencia}
	\label{sug:3.4.10}
	Para el Problema \ref{prob:3.4.10}, exprese la masa de probabilidad del suma como $P(Z = m) = \sum_{j=0}^m P(X = j) P(Y = m-j) = \sum_{j=0}^m \comb{n_1}{j}\comb{n_2}{m-j} p^m q^{(n_1+n_2)-m}$. Aplique la Identidad de Vandermonde: $\sum_{j=0}^m \comb{n_1}{j}\comb{n_2}{m-j} = \comb{n_1+n_2}{m}$.
\end{sugerencia}

\subsection*{Soluciones de los problemas --- Sección 03.04}

\begin{solucion}
	\label{sol:3.4.1}
	Para el Problema \ref{prob:3.4.1}:
	\begin{enumerate}
		\item Para el componente electrónico, el soporte discreto es $\mathcal{S}_X = \set{0, 1}$. Evaluar la forma paramétrica compacta arroja:
		\begin{align*}
			f_X(1) &= p^1 (1-p)^{1-1} = p = 0.85, \\
			f_X(0) &= p^0 (1-p)^{1-0} = 1-p = 0.15.
		\end{align*}
		La suma sobre el soporte es $\sum_{x \in \set{0,1}} f_X(x) = p + (1-p) = 1$, confirmando su normalización exacta.
		\item Calculamos el primer momento bruto (esperanza matemática):
		\begin{align*}
			\mathbb{E}[X] &= \sum_{x \in \set{0,1}} x \cdot f_X(x) = 0 \cdot (1-p) + 1 \cdot p = p.
		\end{align*}
		Para el segundo momento bruto, al ser $x^2 = x$ para todo $x \in \set{0, 1}$:
		\begin{align*}
			\mathbb{E}[X^2] &= \sum_{x \in \set{0,1}} x^2 \cdot f_X(x) = 0^2 \cdot (1-p) + 1^2 \cdot p = p.
		\end{align*}
		Por la identidad de König-Huygens, la varianza es:
		\begin{align*}
			\text{Var}(X) &= \mathbb{E}[X^2] - (\mathbb{E}[X])^2 = p - p^2 = p(1-p).
		\end{align*}
		Sustituyendo $p = 0.85$, obtenemos $\mathbb{E}[X] = 0.85$ y $\text{Var}(X) = 0.85 \times 0.15 = 0.1275$.
	\end{enumerate}
\end{solucion}

\begin{solucion}
	\label{sol:3.4.2}
	Para el Problema \ref{prob:3.4.2}:
	\begin{enumerate}
		\item La variable aleatoria que cuenta éxitos independientes entre un número fijo de ensayos sigue una distribución binomial: $X \sim \text{Bin}(n=8, p=0.90)$. La probabilidad exacta de obtener $7$ sensores aptos es:
		\begin{align*}
			P(X = 7) &= \comb{8}{7} (0.90)^7 (0.10)^{8-7} \\
			&= 8 \times 0.4782969 \times 0.10 \approx 0.3826.
		\end{align*}
		\item Que haya a lo sumo $1$ sensor no apto significa que el número de sensores defectuosos ($8-X$) es $\le 1$, lo cual ocurre si y solo si los sensores aptos son exactamente $7$ u $8$:
		\begin{align*}
			P(X \ge 7) &= P(X = 7) + P(X = 8) \\
			&= 0.38263752 + \comb{8}{8} (0.90)^8 (0.10)^0 \\
			&= 0.38263752 + 1 \times 0.43046721 \approx 0.8131.
		\end{align*}
		Existe un $81.31\%$ de probabilidad de que el lote inspeccionado tenga a lo sumo un fallo.
	\end{enumerate}
\end{solucion}

\begin{solucion}
	\label{sol:3.4.3}
	Para el Problema \ref{prob:3.4.3}:
	\begin{enumerate}
		\item El espacio muestral completo del experimento consta de $2^n$ secuencias ordenadas de ceros y unos de la forma $\omega = (x_1, x_2, \dots, x_n)$. Como cada variable es independiente, la probabilidad de cualquier secuencia particular que contenga exactamente $k$ unos (éxitos) y $n-k$ ceros (fracasos) está dada por el producto de sus probabilidades marginales:
		\begin{align*}
			P(\omega) &= \prod_{i=1}^n P(X_i = x_i) = p^k q^{n-k}.
		\end{align*}
		El número total de secuencias en $\Omega$ que contienen exactamente $k$ éxitos es el problema combinatorio de elegir las $k$ posiciones de los éxitos entre los $n$ lugares disponibles, lo cual es el coeficiente combinatorio $\comb{n}{k} = \dfrac{n!}{k!(n-k)!}$.
		\item Dado que las secuencias que producen $Y=k$ son eventos mutuamente excluyentes y equiprobables entre sí con probabilidad $p^k q^{n-k}$, sumando sobre las $\comb{n}{k}$ configuraciones obtenemos la función de masa de probabilidad binomial:
		\begin{align*}
			P(Y = k) = \comb{n}{k} p^k q^{n-k}, \quad k \in \set{0, 1, \dots, n}.
		\end{align*}
		Su suma sobre todo el soporte $\mathcal{S}_Y$ es:
		\begin{align*}
			\sum_{k=0}^n P(Y = k) = \sum_{k=0}^n \comb{n}{k} p^k q^{n-k} = (p + q)^n = 1^n = 1,
		\end{align*}
		lo cual confirma que es una función de masa de probabilidad matemáticamente completa.
	\end{enumerate}
\end{solucion}

\begin{solucion}
	\label{sol:3.4.4}
	Para el Problema \ref{prob:3.4.4}:
	\begin{enumerate}
		\item Con $X \sim \text{Bin}(12, 0.15)$, calculamos las probabilidades acumuladas hasta $k=2$:
		\begin{align*}
			P(X = 0) &= \comb{12}{0} (0.15)^0 (0.85)^{12} = (0.85)^{12} \approx 0.14224, \\
			P(X = 1) &= \comb{12}{1} (0.15)^1 (0.85)^{11} = 12 \times 0.15 \times 0.16734 \approx 0.30121, \\
			P(X = 2) &= \comb{12}{2} (0.15)^2 (0.85)^{10} = 66 \times 0.0225 \times 0.19687 \approx 0.29236.
		\end{align*}
		Sumando las tres probabilidades elementales:
		\begin{align*}
			P(X \le 2) = 0.14224 + 0.30121 + 0.29236 = 0.7358.
		\end{align*}
		\item Para la probabilidad condicional de al menos $4$ detecciones dado al menos una, calculamos primero $P(X = 3)$:
		\begin{align*}
			P(X = 3) &= \comb{12}{3} (0.15)^3 (0.85)^9 = 220 \times 0.003375 \times 0.23162 \approx 0.17197.
		\end{align*}
		Por lo tanto, la probabilidad de tener $3$ o menos detecciones es $P(X \le 3) = 0.73581 + 0.17197 = 0.90778$. La probabilidad condicional requerida es:
		\begin{align*}
			P(X \ge 4 \mid X \ge 1) &= \dfrac{P(X \ge 4 \cap X \ge 1)}{P(X \ge 1)} = \dfrac{P(X \ge 4)}{P(X \ge 1)} = \dfrac{1 - P(X \le 3)}{1 - P(X = 0)} \\
			&= \dfrac{1 - 0.90778}{1 - 0.14224} = \dfrac{0.09222}{0.85776} \approx 0.1075.
		\end{align*}
	\end{enumerate}
\end{solucion}

\begin{solucion}
	\label{sol:3.4.5}
	Para el Problema \ref{prob:3.4.5}:
	\begin{enumerate}
		\item Sea $X \sim \text{Bin}(n, p=0.08)$ el número de secuencias detectadas. Queremos garantizar que:
		\begin{align*}
			P(X \ge 1) \ge 0.99.
		\end{align*}
		Utilizando el evento complementario, $P(X \ge 1) = 1 - P(X = 0) = 1 - (1-p)^n = 1 - (0.92)^n$. Sustituyendo en la desigualdad:
		\begin{align*}
			1 - (0.92)^n \ge 0.99 \implies (0.92)^n \le 0.01.
		\end{align*}
		Aplicando logaritmo natural (y dado que $\ln(0.92) < 0$, el signo de la desigualdad se invierte):
		\begin{align*}
			n \cdot \ln(0.92) \le \ln(0.01) \implies n \ge \dfrac{\ln(0.01)}{\ln(0.92)}.
		\end{align*}
		\item Evaluando los logaritmos numéricamente con alta precisión:
		\begin{align*}
			\ln(0.01) \approx -4.605170, \quad \ln(0.92) \approx -0.083382.
		\end{align*}
		El cociente exacto es:
		\begin{align*}
			n \ge \dfrac{-4.605170}{-0.083382} \approx 55.2299.
		\end{align*}
		Como el tamaño de muestra $n$ debe ser un número entero positivo, aplicamos la función techo:
		\begin{align*}
			n^* = \lceil 55.2299 \rceil = 56.
		\end{align*}
		Se requieren exactamente $56$ muestras independientes para asegurar el nivel de detección deseado.
	\end{enumerate}
\end{solucion}

\begin{solucion}
	\label{sol:3.4.6}
	Para el Problema \ref{prob:3.4.6}:
	\begin{enumerate}
		\item Para $X \sim \text{Bin}(n=20, p=0.05)$, calculamos primero la media $\mathbb{E}[X]$ y la varianza $\text{Var}(X)$:
		\begin{align*}
			\mathbb{E}[X] &= np = 20 \times 0.05 = 1.0, \\
			\text{Var}(X) &= np(1-p) = 20 \times 0.05 \times 0.95 = 0.95.
		\end{align*}
		El segundo momento bruto es $\mathbb{E}[X^2] = \text{Var}(X) + (\mathbb{E}[X])^2 = 0.95 + (1.0)^2 = 1.95$. Aplicando linealidad al costo de garantía $C(X)$:
		\begin{align*}
			\mathbb{E}[C(X)] &= 500 \mathbb{E}[X^2] + 150 \mathbb{E}[X] + 1000 \\
			&= 500(1.95) + 150(1.0) + 1000 = 975 + 150 + 1000 = \$2,125.
		\end{align*}
		\item La desviación estándar del número de componentes defectuosos en el lote es:
		\begin{align*}
			\sigma_X = \sqrt{\text{Var}(X)} = \sqrt{0.95} \approx 0.9747.
		\end{align*}
		Si el costo fuera puramente lineal en la media ($150 \mathbb{E}[X] + 1000$), el costo esperado sería de $\$1,150$. Sin embargo, la no linealidad cuadrática ($500 X^2$) penaliza fuertemente la variabilidad inherente del lote: aun cuando el número promedio de fallas es solo una unidad ($\mu=1$), la dispersión intrínseca agrega un costo esperado de $500 \times 1.95 = \$975$, lo cual representa el $45.88\%$ del costo global.
	\end{enumerate}
\end{solucion}

\begin{solucion}
	\label{sol:3.4.7}
	Para el Problema \ref{prob:3.4.7}:
	\begin{enumerate}
		\item El Teorema del Binomio de Newton establece que para cualquier par de números reales $a, b \in \mathbb{R}$ y entero positivo $n \in \mathbb{Z}^+$:
		\begin{align*}
			(a + b)^n = \sum_{k=0}^n \comb{n}{k} a^k b^{n-k}.
		\end{align*}
		Sustituyendo directamente $a = p$ y $b = q$, obtenemos el lado derecho idéntico a la suma total de probabilidades de la función de masa binomial:
		\begin{align*}
			\sum_{k=0}^n \comb{n}{k} p^k q^{n-k} = (p + q)^n.
		\end{align*}
		Como por definición de probabilidad de fracaso se cumple que $p + q = 1$, la expresión se reduce a $1^n = 1$, demostrando formalmente la normalización absoluta.
		\item Tratando a $p$ y $q$ como variables independientes en la identidad general $(p+q)^n = \sum_{k=0}^n \comb{n}{k} p^k q^{n-k}$ y aplicando el operador diferencial respecto a $p$:
		\begin{align*}
			\dfrac{\partial}{\partial p} (p+q)^n &= \sum_{k=0}^n \comb{n}{k} \dfrac{\partial}{\partial p} (p^k q^{n-k}) \\
			n(p+q)^{n-1} &= \sum_{k=1}^n k \comb{n}{k} p^{k-1} q^{n-k}.
		\end{align*}
		Multiplicamos ambos lados de la ecuación por el parámetro $p$:
		\begin{align*}
			np(p+q)^{n-1} = \sum_{k=1}^n k \comb{n}{k} p^k q^{n-k} = \sum_{k=0}^n k \comb{n}{k} p^k q^{n-k}.
		\end{align*}
		El lado derecho de esta ecuación coincide exactamente con la definición analítica del valor esperado $\mathbb{E}[X] = \sum_{k=0}^n k P(X = k)$. Al sustituir la condición geométrica $q = 1-p \implies p+q=1$ en el lado izquierdo, concluimos la demostración:
		\begin{align*}
			\mathbb{E}[X] = np(1)^{n-1} = np.
		\end{align*}
	\end{enumerate}
\end{solucion}

\begin{solucion}
	\label{sol:3.4.8}
	Para el Problema \ref{prob:3.4.8}:
	\begin{enumerate}
		\item Evaluamos la definición formal de segundo momento factorial $\mathbb{E}[X(X-1)]$:
		\begin{align*}
			\mathbb{E}[X(X-1)] &= \sum_{k=0}^n k(k-1) P(X = k) = \sum_{k=2}^n k(k-1) \comb{n}{k} p^k q^{n-k}.
		\end{align*}
		Desarrollamos el coeficiente combinatorio y cancelamos los factores algebraicos:
		\begin{align*}
			k(k-1) \comb{n}{k} &= k(k-1) \dfrac{n!}{k!(n-k)!} = \dfrac{n!}{(k-2)!(n-k)!} \\
			&= n(n-1) \dfrac{(n-2)!}{(k-2)!((n-2)-(k-2))!} = n(n-1) \comb{n-2}{k-2}.
		\end{align*}
		Sustituyendo este desarrollo combinatorio en la sumatoria y factorizando $n(n-1)p^2$:
		\begin{align*}
			\mathbb{E}[X(X-1)] &= \sum_{k=2}^n n(n-1)\comb{n-2}{k-2} p^2 p^{k-2} q^{(n-2)-(k-2)} \\
			&= n(n-1)p^2 \sum_{j=0}^{n-2} \comb{n-2}{j} p^j q^{(n-2)-j}, \quad \text{donde } j = k-2.
		\end{align*}
		Por el Teorema del Binomio demostrado previamente, la sumatoria interna vale exactamente $(p+q)^{n-2} = 1$. Por lo tanto:
		\begin{align*}
			\mathbb{E}[X(X-1)] = n(n-1)p^2.
		\end{align*}
		\item Desarrollamos el momento factorial relacionándolo con la varianza de la distribución:
		\begin{align*}
			\mathbb{E}[X(X-1)] &= \mathbb{E}[X^2 - X] = \mathbb{E}[X^2] - \mathbb{E}[X] \\
			\implies \mathbb{E}[X^2] &= \mathbb{E}[X(X-1)] + \mathbb{E}[X].
		\end{align*}
		Sustituyendo $\mathbb{E}[X(X-1)] = n(n-1)p^2$ y $\mathbb{E}[X] = np$, obtenemos el segundo momento bruto:
		\begin{align*}
			\mathbb{E}[X^2] = n(n-1)p^2 + np = n^2 p^2 - np^2 + np.
		\end{align*}
		Finalmente, calculamos la varianza operacional:
		\begin{align*}
			\text{Var}(X) &= \mathbb{E}[X^2] - (\mathbb{E}[X])^2 = (n^2 p^2 - np^2 + np) - (np)^2 \\
			&= np - np^2 = np(1-p),
		\end{align*}
		concluyendo la deducción de la varianza binomial de manera exacta.
	\end{enumerate}
\end{solucion}

\begin{solucion}
	\label{sol:3.4.9}
	Para el Problema \ref{prob:3.4.9}:
	\begin{enumerate}
		\item Para determinar dónde alcanza su máximo la función de masa binomial, investigamos cuándo la probabilidad crece de un término al siguiente ($P(X=k) \ge P(X=k-1)$). Evaluamos el cociente adyacente $R(k)$:
		\begin{align*}
			R(k) = \dfrac{P(X=k)}{P(X=k-1)} &= \dfrac{\comb{n}{k} p^k q^{n-k}}{\comb{n}{k-1} p^{k-1} q^{n-k+1}} = \dfrac{\dfrac{n!}{k!(n-k)!}}{\dfrac{n!}{(k-1)!(n-k+1)!}} \cdot \dfrac{p}{q} \\
			&= \dfrac{(k-1)!(n-k+1)!}{k!(n-k)!} \cdot \dfrac{p}{q} = \dfrac{n-k+1}{k} \cdot \dfrac{p}{1-p}.
		\end{align*}
		\item La probabilidad binomial es creciente en el paso $k$ si y solo si $R(k) \ge 1$:
		\begin{align*}
			\dfrac{n-k+1}{k} \cdot \dfrac{p}{1-p} \ge 1 &\implies (n-k+1)p \ge k(1-p) \\
			&\implies np - kp + p \ge k - kp \\
			&\implies (n+1)p \ge k.
		\end{align*}
		Por lo tanto, las probabilidades son estrictamente crecientes en los enteros $k$ tales que $k \le (n+1)p$, y estrictamente decrecientes para los enteros $k > (n+1)p$. Cuando la cantidad $(n+1)p$ no es un número entero, existe un único número entero $k^*$ que satisface la condición de máximo, el cual se obtiene tomando la parte entera inferior:
		\begin{align*}
			k^* = \lfloor (n+1)p \rfloor.
		\end{align*}
		\item Si el valor admisible $(n+1)p = m$ resulta ser un entero positivo exacto, entonces en el punto de transición $k = m$ se tiene una igualdad en el cociente adyacente:
		\begin{align*}
			R(m) = 1 \implies P(X = m) = P(X = m-1).
		\end{align*}
		En este caso extraordinario la distribución es exactamente **bimodal**, alcanzando su máximo valor idéntico en dos puntos enteros contiguos del soporte: $k^* = m-1$ y $k^* = m$.
	\end{enumerate}
\end{solucion}

\begin{solucion}
	\label{sol:3.4.10}
	Para el Problema \ref{prob:3.4.10}:
	\begin{enumerate}
		\item Al ser $X$ y $Y$ independientes con el mismo éxito $p$, la probabilidad de que la suma tome un valor entero $Z = X+Y = m$ para $m \in \set{0, 1, \dots, n_1+n_2}$ se obtiene mediante la fórmula de convolución discreta:
		\begin{align*}
			P(Z = m) &= \sum_{j=0}^m P(X = j \cap Y = m-j) = \sum_{j=0}^m P(X = j)P(Y = m-j) \\
			&= \sum_{j=0}^m \left[ \comb{n_1}{j} p^j q^{n_1-j} \right] \left[ \comb{n_2}{m-j} p^{m-j} q^{n_2-(m-j)} \right] \\
			&= p^m q^{(n_1+n_2)-m} \sum_{j=0}^m \comb{n_1}{j} \comb{n_2}{m-j}.
		\end{align*}
		La sumatoria combinatoria remanente se resuelve mediante la célebre **Identidad combinatoria de Vandermonde**, la cual establece que al seleccionar un subconjunto de tamaño $m$ de un grupo combinado de $n_1 + n_2$ elementos divididos en dos subgrupos disjuntos:
		\begin{align*}
			\sum_{j=0}^m \comb{n_1}{j} \comb{n_2}{m-j} = \comb{n_1+n_2}{m}.
		\end{align*}
		Sustituyendo la Identidad de Vandermonde en la convolución, obtenemos de forma cerrada:
		\begin{align*}
			P(Z = m) = \comb{n_1+n_2}{m} p^m q^{(n_1+n_2)-m},
		\end{align*}
		lo que demuestra que $Z \sim \text{Bin}(n_1+n_2, p)$.
		\item Verificamos la aditividad de momentos sobre la suma binomial:
		\begin{align*}
			\mathbb{E}[Z] &= (n_1 + n_2)p = n_1 p + n_2 p = \mathbb{E}[X] + \mathbb{E}[Y], \\
			\text{Var}(Z) &= (n_1 + n_2)p(1-p) = n_1 p(1-p) + n_2 p(1-p) = \text{Var}(X) + \text{Var}(Y).
		\end{align*}
		El teorema combinatorio y aditivo queda rigurosamente demostrado.
	\end{enumerate}
\end{solucion}

% ==============================================================================
% SECCIÓN 03.05: DISTRIBUCIONES GEOMÉTRICA Y BINOMIAL NEGATIVA
% ==============================================================================

\subsection*{Enunciados de los problemas --- Sección 03.05}

\subsubsection*{Nivel Fundamental}

\begin{problema}
	\label{prob:3.5.1}
	Considere una secuencia infinita de ensayos de Bernoulli independientes con probabilidad constante de éxito $p \in (0, 1)$ en cada intento y probabilidad de fracaso $q = 1-p$. Sea $X$ la variable aleatoria discreta que denota el número total de ensayos requeridos hasta obtener el **primer éxito** (incluyendo el ensayo exitoso).
	\begin{enumerate}
		\item Especifique la función de masa de probabilidad $f_X(k) = P(X = k)$ y determine su soporte $\mathcal{S}_X$.
		\item Demuestre analíticamente que la suma sobre todo el soporte es igual a la unidad ($\sum_{k=1}^{\infty} f_X(k) = 1$) mediante series geométricas.
		\item Calcule la probabilidad de cola superior $P(X > k)$ para un entero $k \ge 1$ y obtenga una expresión cerrada para la función de distribución acumulada $F_X(k) = P(X \le k)$.
	\end{enumerate}
\end{problema}

\begin{problema}
	\label{prob:3.5.2}
	En un servidor de base de datos de alta concurrencia, las transacciones independientes experimentan fallas de bloqueo por concurrencia con probabilidad constante $p = 0.04$. Sea $X$ el número de intentos necesarios para que una transacción compleja logre ejecutarse con éxito por primera vez.
	\begin{enumerate}
		\item Calcule la probabilidad exacta de que la transacción requiera a lo sumo $3$ intentos para ejecutarse exitosamente.
		\item Si la transacción ya ha fallado en sus primeros $5$ intentos, ¿cuál es la probabilidad condicional de que requiera más de $8$ intentos en total para tener éxito ($P(X > 8 \mid X > 5)$)?
		\item Compare el resultado condicional anterior con la probabilidad no condicional de requerir más de $3$ intentos ($P(X > 3)$) e interprete el resultado en términos de la propiedad intrínseca del proceso.
	\end{enumerate}
\end{problema}

\begin{problema}
	\label{prob:3.5.3}
	Una firma de auditoría contable inspecciona transacciones corporativas de manera secuencial e independiente. Históricamente, se sabe que una fracción $p = 0.10$ de las transacciones presenta alguna irregularidad fiscal. Los auditores deciden continuar inspeccionando hasta encontrar **exactamente $r = 3$ transacciones irregulares**. Sea $X$ el número total de transacciones inspeccionadas.
	\begin{enumerate}
		\item Argumente por qué la probabilidad de que la tercera irregularidad sea detectada exactamente en la vigésima inspección ($X = 20$) adopta la forma combinatoria $P(X = 20) = \comb{19}{2} (0.10)^3 (0.90)^{17}$.
		\item Escriba la función de masa general para $X \sim \text{BN}(r, p)$ sobre su soporte $\mathcal{S}_X = \set{r, r+1, \dots}$.
	\end{enumerate}
\end{problema}

\subsubsection*{Nivel Operativo}

\begin{problema}
	\label{prob:3.5.4}
	En una línea de soldadura láser automatizada para componentes de automoción, los defectos críticos ocurren de manera independiente con probabilidad $p = 0.02$ por pieza. Por protocolo de seguridad industrial, la línea se detiene para un recalibrado mayor en el instante exacto en que se detecta la **cuarta pieza defectuosa** ($r = 4$). Sea $X$ el número total de piezas producidas hasta el paro de la línea.
	\begin{enumerate}
		\item Calcule el número esperado de piezas producidas hasta el recalibrado ($\mathbb{E}[X]$) y la desviación estándar de la producción ($\sigma_X$).
		\item Calcule la probabilidad de que la línea de soldadura logre producir al menos $100$ piezas antes de ser detenida por el cuarto defecto ($P(X \ge 100)$).
	\end{enumerate}
\end{problema}

\begin{problema}
	\label{prob:3.5.5}
	Un equipo de biotecnología realiza microinyecciones en ovocitos de ratón para desarrollar un modelo transgénico. La probabilidad de éxito de cada microinyección es un parámetro desconocido $p \in (0, 1)$. El equipo planea inyectar ovocitos uno tras otro hasta obtener exactamente $r = 5$ embriones viables. Suponga que en el experimento real, el quinto embrión viable se obtuvo exactamente en la microinyección número $25$ ($X = 25$).
	\begin{enumerate}
		\item Escriba la función de verosimilitud $L(p \mid X = 25)$ basada en la distribución Binomial Negativa.
		\item Calcule el estimador de máxima verosimilitud $\hat{p}_{\text{MLE}}$ para la probabilidad de viabilidad e interprete su relación intuitiva con la proporción observada de éxitos.
	\end{enumerate}
\end{problema}

\begin{problema}
	\label{prob:3.5.6}
	Un ingeniero en telecomunicaciones analiza el número de reintentos de conexión por usuario en una hora punta del servidor central. Al ajustar el número de fallas previas al primer éxito con un modelo discreto, observa que la media muestral es $\bar{x} = 3.2$ fallas y la varianza muestral es $s^2 = 13.44$.
	\begin{enumerate}
		\item Explique teóricamente por qué un modelo de distribución de Poisson pura resulta estructuralmente inadecuado para modelar este tráfico de reintentos.
		\item Si se modela el número de fallas $Y = X - r \sim \text{BN}(r, p)$ (donde $Y \in \set{0, 1, \dots}$ tiene media $\mu_Y = \frac{rq}{p}$ y varianza $\sigma_Y^2 = \frac{rq}{p^2}$), utilice el método de momentos igualando $\mu_Y = 3.2$ y $\sigma_Y^2 = 13.44$ para estimar los parámetros poblacionales $\hat{p}$ y $\hat{r}$.
	\end{enumerate}
\end{problema}

\subsubsection*{Nivel Analítico}

\begin{problema}
	\label{prob:3.5.7}
	Sea un experimento secuencial donde $X \sim \text{BN}(r, p)$ cuenta el número total de ensayos de Bernoulli necesarios para acumular $r$ éxitos independientes.
	\begin{enumerate}
		\item Demuestre rigurosamente por descomposición estocástica que $X$ puede expresarse como la suma de $r$ variables aleatorias geométricas independientes e idénticamente distribuidas ($X = \sum_{i=1}^r Y_i$, donde $Y_i \sim \text{Geom}(p)$).
		\item Demuestre por diferenciación de series de potencias que para una única variable $Y_1 \sim \text{Geom}(p)$, la función generadora de momentos es $M_{Y_1}(t) = \dfrac{p e^t}{1 - q e^t}$ para $t < -\ln(q)$.
		\item Utilizando la propiedad de independencia, deduzca la función generadora de momentos $M_X(t)$ de la distribución Binomial Negativa.
	\end{enumerate}
\end{problema}

\begin{problema}
	\label{prob:3.5.8}
	Utilizando las propiedades analíticas del operador diferencial sobre series formales de potencias o la función generadora de momentos $M_X(t) = \left(\dfrac{p e^t}{1 - q e^t}\right)^r$ obtenida en el Problema \ref{prob:3.5.7}:
	\begin{enumerate}
		\item Demuestre de manera exacta que la esperanza matemática del número total de ensayos en una Binomial Negativa es $\mathbb{E}[X] = \dfrac{r}{p}$.
		\item Demuestre algebraicamente que la varianza está dada por $\text{Var}(X) = \dfrac{r q}{p^2} = \dfrac{r(1-p)}{p^2}$.
	\end{enumerate}
\end{problema}

\subsubsection*{Nivel Desafiante}

\begin{problema}
	\label{prob:3.5.9}
	Se afirma que la distribución geométrica es la **única distribución discreta con soporte en los enteros positivos $\mathbb{Z}^+$** que posee la propiedad de pérdida de memoria. Demuestre formalmente este teorema de caracterización:
	\begin{align*}
		P(X > m+n \mid X > m) = P(X > n) \quad \forall m, n \in \mathbb{Z}^+ \iff X \sim \text{Geom}(p).
	\end{align*}
\end{problema}

\begin{problema}
	\label{prob:3.5.10}
	Considere la parametrización del número de fracasos observados antes de obtener $r$ éxitos: $Y = X - r \sim \text{BN}(r, p)$ con soporte $Y \in \set{0, 1, 2, \dots}$ y función de masa $P(Y = y) = \comb{y+r-1}{y} p^r (1-p)^y$. Suponga que el requisito de éxitos crece ilimitadamente ($r \to \infty$) mientras que la probabilidad de éxito de cada ensayo individual tiende a la certeza ($p \to 1$), de tal forma que el número medio esperado de fracasos se mantiene constante e igual a una tasa fija $\lambda > 0$. Esto es, $(1-p)r = \lambda \implies 1-p = \frac{\lambda}{r}$.
	\begin{enumerate}
		\item Demuestre analíticamente mediante límites combinatorios y exponenciales que para cualquier entero no negativo fijo $y \in \set{0, 1, \dots}$:
		\begin{align*}
			\lim_{r \to \infty} P(Y = y) = \dfrac{\lambda^y e^{-\lambda}}{y!}.
		\end{align*}
		\item Interprete físicamente este límite en el contexto del modelado de fenómenos raros y la transición del esquema combinatorio de Pascal al proceso puntual de Poisson.
	\end{enumerate}
\end{problema}

% ------------------------------------------------------------------------------
% SUGERENCIAS --- SECCIÓN 03.05
% ------------------------------------------------------------------------------

\subsection*{Sugerencias de los problemas --- Sección 03.05}

\begin{sugerencia}
	\label{sug:3.5.1}
	Para la suma en el soporte, recuerde la serie geométrica infinita $\sum_{k=0}^{\infty} x^k = \frac{1}{1-x}$ válida cuando $|x| < 1$. Para la cola superior $P(X > k)$, sume la serie desde $j = k+1$ hasta $\infty$ factorizando la potencia común $q^k$.
\end{sugerencia}

\begin{sugerencia}
	\label{sug:3.5.2}
	Exprese la probabilidad condicional utilizando la definición básica: $P(A \mid B) = \frac{P(A \cap B)}{P(B)}$. Note que el evento $\set{X > m+n} \cap \set{X > m}$ es simplemente $\set{X > m+n}$. Sustituya $P(X > j) = q^j$ deducido en el Problema \ref{prob:3.5.1}.
\end{sugerencia}

\begin{sugerencia}
	\label{sug:3.5.3}
	Para que el $r$-ésimo éxito ocurra en el intento exacto $k$, deben cumplirse dos condiciones simultáneas e independientes: (1) haber acumulado exactamente $r-1$ éxitos en los primeros $k-1$ intentos, y (2) que el intento $k$ sea un éxito. Multiplique ambas probabilidades combinando la Binomial y Bernoulli.
\end{sugerencia}

\begin{sugerencia}
	\label{sug:3.5.4}
	Aplique las fórmulas teóricas de momentos de la Binomial Negativa: $\mathbb{E}[X] = \frac{r}{p}$ y $\text{Var}(X) = \frac{rq}{p^2}$. Para la probabilidad $P(X \ge 100)$, reconozca la equivalencia con el conteo binomial: necesitar al menos $100$ ensayos para lograr $4$ defectos equivale a observar a lo sumo $3$ defectos en los primeros $99$ ensayos fabricados.
\end{sugerencia}

\begin{sugerencia}
	\label{sug:3.5.5}
	Tome el logaritmo natural de la verosimilitud $\ln L(p) = \ln \comb{k-1}{r-1} + r \ln(p) + (k-r)\ln(1-p)$. Derive respecto a $p$, iguale a cero y despeje $p$ en términos de $r$ y $k$.
\end{sugerencia}

\begin{sugerencia}
	\label{sug:3.5.6}
	Recuerde que en una distribución de Poisson pura $\mu = \sigma^2$. Para el sistema de ecuaciones con $\text{BN}$, observe el cociente $\frac{\sigma_Y^2}{\mu_Y} = \frac{rq/p^2}{rq/p} = \frac{1}{p} = 1 + \frac{q}{p}$. Despeje $p$ del cociente varianza-media y luego sustitúyalo en $\mu_Y$ para obtener $r$.
\end{sugerencia}

\begin{sugerencia}
	\label{sug:3.5.7}
	El tiempo hasta el $r$-ésimo éxito es la suma del tiempo hasta el 1er éxito ($Y_1$), más el tiempo adicional entre el 1er y el 2do éxito ($Y_2$), y así sucesivamente hasta $Y_r$. Para la FGM, sume la serie $\sum_{k=1}^{\infty} e^{tk} q^{k-1} p$ reescribiéndola como $p e^t \sum_{j=0}^{\infty} (q e^t)^j$.
\end{sugerencia}

\begin{sugerencia}
	\label{sug:3.5.8}
	Puede derivar $M_X(t) = \left(\frac{p e^t}{1-qe^t}\right)^r$ una y dos veces respecto a $t$ y evaluar en $t=0$, o de manera más ágil, aprovechar la independencia y linealidad de momentos proveniente de $X = \sum_{i=1}^r Y_i$ usando los momentos de la geométrica elemental.
\end{sugerencia}

\begin{sugerencia}
	\label{sug:3.5.9}
	Defina la función de cola discreta $G(k) = P(X > k)$. La hipótesis de pérdida de memoria se traduce en la ecuación funcional discreta $G(m+n) = G(m)G(n)$ para todo $m, n \ge 1$. Por inducción sobre $n$, demuestre que $G(n) = [G(1)]^n$ y relacione $G(1)$ con la probabilidad de fracaso $q = 1-p$.
\end{sugerencia}

\begin{sugerencia}
	\label{sug:3.5.10}
	Sustituya $p = 1 - \frac{\lambda}{r}$ en la PMF binomial negativa. Expanda el coeficiente combinatorio $\comb{y+r-1}{y} = \frac{(r+y-1)(r+y-2)\cdots r}{y!}$ y agrupe los $y$ términos que crecen con $r$. Utilice el límite exponencial fundamental $\lim_{r\to\infty} \left(1 - \frac{\lambda}{r}\right)^r = e^{-\lambda}$.
\end{sugerencia}

% ------------------------------------------------------------------------------
% SOLUCIONES --- SECCIÓN 03.05
% ------------------------------------------------------------------------------

\subsection*{Soluciones de los problemas --- Sección 03.05}

\begin{solucion}
	\label{sol:3.5.1}
	Para el Problema \ref{prob:3.5.1}:
	\begin{enumerate}
		\item Para que el primer éxito ocurra exactamente en el ensayo $k \ge 1$, los primeros $k-1$ ensayos debieron resultar en fracasos consecutivos (probabilidad $q^{k-1}$ por independencia mutua) y el ensayo $k$ debe ser un éxito (probabilidad $p$). Por lo tanto:
		\begin{align*}
			f_X(k) = P(X = k) = q^{k-1}p = (1-p)^{k-1}p, \quad \text{para } k \in \mathcal{S}_X = \set{1, 2, 3, \dots}.
		\end{align*}
		\item Sumamos la probabilidad sobre todo el soporte de números enteros positivos:
		\begin{align*}
			\sum_{k=1}^{\infty} f_X(k) &= \sum_{k=1}^{\infty} q^{k-1}p = p \sum_{j=0}^{\infty} q^j, \quad \text{haciendo el cambio de índice } j = k-1.
		\end{align*}
		Como $p \in (0, 1)$, se tiene $q = 1-p \in (0, 1)$. Aplicando la suma exacta de la serie geométrica convergente:
		\begin{align*}
			\sum_{j=0}^{\infty} q^j = \dfrac{1}{1 - q} = \dfrac{1}{1 - (1-p)} = \dfrac{1}{p}.
		\end{align*}
		Sustituyendo en la sumatoria original, obtenemos $p \cdot \dfrac{1}{p} = 1$, lo que demuestra la normalización absoluta de la distribución.
		\item La probabilidad de cola superior $P(X > k)$ representa el evento de que los primeros $k$ intentos sean todos fracasos continuos. Sumando la serie analítica:
		\begin{align*}
			P(X > k) &= \sum_{j=k+1}^{\infty} q^{j-1}p = p q^k \sum_{m=0}^{\infty} q^m = p q^k \left(\dfrac{1}{1-q}\right) = p q^k \left(\dfrac{1}{p}\right) = q^k.
		\end{align*}
		Por complementariedad, la función de distribución acumulada para cualquier entero $k \ge 1$ es:
		\begin{align*}
			F_X(k) = P(X \le k) = 1 - P(X > k) = 1 - q^k = 1 - (1-p)^k.
		\end{align*}
	\end{enumerate}
\end{solucion}

\begin{solucion}
	\label{sol:3.5.2}
	Para el Problema \ref{prob:3.5.2}:
	\begin{enumerate}
		\item Con probabilidad de éxito $p = 1 - 0.04 = 0.96$ y fracaso $q = 0.04$:
		\begin{align*}
			P(X \le 3) = F_X(3) = 1 - q^3 = 1 - (0.04)^3 = 1 - 0.000064 = 0.999936.
		\end{align*}
		Existe una certeza práctica de que la transacción se procesará en máximo $3$ intentos.
		\item Aplicando la probabilidad condicional con los eventos de cola deducidos en el Problema \ref{prob:3.5.1} ($P(X > k) = q^k$):
		\begin{align*}
			P(X > 8 \mid X > 5) &= \dfrac{P(\set{X > 8} \cap \set{X > 5})}{P(X > 5)} = \dfrac{P(X > 8)}{P(X > 5)} \\
			&= \dfrac{q^8}{q^5} = q^{8-5} = q^3 = (0.04)^3 = 0.000064.
		\end{align*}
		\item Observamos que la probabilidad no condicional de requerir más de $3$ intentos desde el inicio del proceso es idéntica:
		\begin{align*}
			P(X > 3) = q^3 = 0.000064 = P(X > 5 + 3 \mid X > 5).
		\end{align*}
		Este resultado ilustra la **propiedad de pérdida de memoria**: el sistema no acumula fatiga ni registra su historial de fallas pasadas. Haber fallado ya $5$ veces no aumenta ni disminuye la probabilidad de requerir $3$ o más intentos adicionales para lograr el éxito.
	\end{enumerate}
\end{solucion}

\begin{solucion}
	\label{sol:3.5.3}
	Para el Problema \ref{prob:3.5.3}:
	\begin{enumerate}
		\item Para que la tercera irregularidad ($r=3$) sea descubierta exactamente en la vigésima auditoría ($X=20$), el evento se divide en dos fases independientes e ineludibles:
		\begin{itemize}
			\item **Fase A (Primeras 19 inspecciones):** Deben haberse detectado exactamente $r-1 = 2$ irregularidades en los primeros $k-1 = 19$ documentos revisados (en cualquier orden). El número de tales combinaciones es el coeficiente binomial $\comb{19}{2}$, y cada una ocurre con probabilidad conjunta $p^2 q^{17}$.
			\item **Fase B (Inspección 20):** El vigésimo documento auditado debe ser obligatoriamente irregular, lo cual tiene probabilidad $p$.
		\end{itemize}
		Por el principio combinatorio multiplicativo:
		\begin{align*}
			P(X = 20) = \left[\comb{19}{2} p^2 q^{17}\right] \times p = \comb{19}{2} p^3 q^{17} = \comb{19}{2} (0.10)^3 (0.90)^{17} \approx 0.0285.
		\end{align*}
		\item Generalizando el argumento anterior para $r$ éxitos en $k$ intentos totales:
		\begin{align*}
			f_X(k) = P(X = k) = \comb{k-1}{r-1} p^r q^{k-r}, \quad \text{para } k \in \mathcal{S}_X = \set{r, r+1, r+2, \dots}.
		\end{align*}
	\end{enumerate}
\end{solucion}

\begin{solucion}
	\label{sol:3.5.4}
	Para el Problema \ref{prob:3.5.4}:
	\begin{enumerate}
		\item Con parámetros de defecto $p = 0.02$, $q = 0.98$ y requisito de paro $r = 4$, aplicamos las expresiones analíticas de momentos:
		\begin{align*}
			\mathbb{E}[X] &= \dfrac{r}{p} = \dfrac{4}{0.02} = 200 \text{ piezas producidas en promedio antes del paro.} \\
			\text{Var}(X) &= \dfrac{rq}{p^2} = \dfrac{4(0.98)}{(0.02)^2} = \dfrac{3.92}{0.0004} = 9,800. \\
			\sigma_X &= \sqrt{\text{Var}(X)} = \sqrt{9,800} \approx 98.9949 \text{ piezas.}
		\end{align*}
		\item El evento de que la línea logre fabricar al menos $100$ piezas antes de ser detenida ($X \ge 100$) equivale al evento dual de que en las primeras $99$ piezas producidas se hayan observado **a lo sumo $3$ defectos**. Sea $W \sim \text{Bin}(n=99, p=0.02)$ el conteo de piezas defectuosas en ese intervalo inicial:
		\begin{align*}
			P(X \ge 100) &= P(W \le 3) = \sum_{j=0}^3 \comb{99}{j} (0.02)^j (0.98)^{99-j} \\
			&= P(W=0) + P(W=1) + P(W=2) + P(W=3).
		\end{align*}
		Calculando los términos individuales algebraicamente:
		\begin{align*}
			P(W=0) &= (0.98)^{99} \approx 0.1353, \\
			P(W=1) &= 99(0.02)(0.98)^{98} \approx 0.2734, \\
			P(W=2) &= 4851(0.0004)(0.98)^{97} \approx 0.2734, \\
			P(W=3) &= 156849(0.000008)(0.98)^{96} \approx 0.1804.
		\end{align*}
		Sumando todas las probabilidades: $P(X \ge 100) \approx 0.1353 + 0.2734 + 0.2734 + 0.1804 = 0.8625$. Existe un $86.25\%$ de probabilidad de superar las $99$ piezas continuas de producción.
	\end{enumerate}
\end{solucion}

\begin{solucion}
	\label{sol:3.5.5}
	Para el Problema \ref{prob:3.5.5}:
	\begin{enumerate}
		\item Para una observación muestral única $X = k = 25$ y $r = 5$, la función de verosimilitud de la probabilidad de viabilidad $p$ es:
		\begin{align*}
			L(p \mid X = 25) = \comb{24}{4} p^5 (1-p)^{20} = 10626 p^5 (1-p)^{20}, \quad \text{para } p \in (0, 1).
		\end{align*}
		\item Tomamos la función de log-verosimilitud $\ell(p) = \ln L(p)$ para simplificar la maximización:
		\begin{align*}
			\ell(p) = \ln(10626) + 5 \ln(p) + 20 \ln(1-p).
		\end{align*}
		Derivamos respecto al parámetro $p$ e igualamos a cero para encontrar la condición de primer orden:
		\begin{align*}
			\dfrac{d \ell(p)}{dp} = \dfrac{5}{p} - \dfrac{20}{1-p} = 0 &\implies 5(1-p) = 20p \\
			&\implies 5 - 5p = 20p \implies 25p = 5 \\
			&\implies \hat{p}_{\text{MLE}} = \dfrac{5}{25} = 0.20.
		\end{align*}
		En el caso general para datos $(r, k)$, el estimador puntual es exactamente $\hat{p}_{\text{MLE}} = \dfrac{r}{k}$, el cual coincide de manera intuitiva y directa con la proporción empírica del número de éxitos deseados sobre el total de intentos realizados.
	\end{enumerate}
\end{solucion}

\begin{solucion}
	\label{sol:3.5.6}
	Para el Problema \ref{prob:3.5.6}:
	\begin{enumerate}
		\item En un proceso puntiforme o de conteo regido por la distribución de Poisson ($\text{Poisson}(\lambda)$), la varianza teórica es idéntica a su esperanza ($\sigma^2 = \mu = \lambda$). Sin embargo, en el tráfico de red de este servidor observamos un cociente de dispersión empírico:
		\begin{align*}
			\text{Índice de Dispersión} = \dfrac{s^2}{\bar{x}} = \dfrac{13.44}{3.2} = 4.2 \gg 1.
		\end{align*}
		La variabilidad empírica es más de $4$ veces superior a la media. Esta **sobredispersión intrínseca** viola frontalmente el supuesto de equidispersión de Poisson, el cual conduciría a subestimar drásticamente las colas y a dimensionar de manera insegura los búferes de conexión.
		\item Utilizando la parametrización en función del conteo de fracasos $Y = X - r \sim \text{BN}(r, p)$, igualamos los momentos teóricos con los muestrales ($\mu_Y = 3.2$ y $\sigma_Y^2 = 13.44$):
		\begin{align*}
			\mu_Y &= \dfrac{r(1-p)}{p} = 3.2, \qquad \sigma_Y^2 = \dfrac{r(1-p)}{p^2} = 13.44.
		\end{align*}
		Tomamos el cociente exacto entre la ecuación de varianza y la ecuación de media:
		\begin{align*}
			\dfrac{\sigma_Y^2}{\mu_Y} = \dfrac{\dfrac{r(1-p)}{p^2}}{\dfrac{r(1-p)}{p}} = \dfrac{1}{p} \implies \dfrac{13.44}{3.2} = \dfrac{1}{\hat{p}} \implies 4.2 = \dfrac{1}{\hat{p}} \implies \hat{p} = \dfrac{1}{4.2} = \dfrac{5}{21} \approx 0.2381.
		\end{align*}
		Sustituyendo $\hat{p} = \frac{5}{21}$ en la ecuación de la media muestral para despejar $\hat{r}$:
		\begin{align*}
			\dfrac{\hat{r}\left(1 - \frac{5}{21}\right)}{\frac{5}{21}} = 3.2 &\implies \hat{r}\left(\dfrac{16/21}{5/21}\right) = 3.2 \implies \hat{r}\left(\dfrac{16}{5}\right) = 3.2 \\
			&\implies \hat{r} = 3.2 \left(\dfrac{5}{16}\right) = \dfrac{16.0}{16} = 1.0.
		\end{align*}
		Los estimadores por el método de momentos son $\hat{r} = 1$ y $\hat{p} \approx 0.2381$. Al ser $r=1$, el modelo que mejor absorbe la sobredispersión de estos datos resulta ser justamente una **Distribución Geométrica** contando reintentos fallidos.
	\end{enumerate}
\end{solucion}

\begin{solucion}
	\label{sol:3.5.7}
	Para el Problema \ref{prob:3.5.7}:
	\begin{enumerate}
		\item Sea $X$ el tiempo total de ensayos hasta el éxito número $r$. Descomponemos el proceso estocástico en incrementos de espera entre éxitos consecutivos:
		\begin{itemize}
			\item Sea $Y_1$ el número de ensayos necesarios para obtener el primer éxito ($\implies Y_1 \sim \text{Geom}(p)$).
			\item Sea $Y_2$ el número adicional de ensayos transcurridos *después* del 1er éxito hasta obtener el 2do éxito. Por la propiedad de pérdida de memoria y la independencia mutua de los ensayos de Bernoulli, la variable del intervalo $Y_2$ es independiente de $Y_1$ y sigue idéntica distribución geométrica $Y_2 \sim \text{Geom}(p)$.
			\item Por inducción finita, para el intervalo $i$-ésimo ($1 \le i \le r$), sea $Y_i$ el número de ensayos adicionales entre el éxito $(i-1)$ y el éxito $i$.
		\end{itemize}
		Sumando todos los intervalos sucesivos, el tiempo total hasta el $r$-ésimo éxito es la suma exacta de las duraciones individuales:
		\begin{align*}
			X = Y_1 + Y_2 + \dots + Y_r = \sum_{i=1}^r Y_i, \quad \text{donde } Y_i \sim \text{Geom}(p) \text{ son i.i.d.}
		\end{align*}
		\item Evaluamos por definición la función generadora de momentos para $Y_1 \sim \text{Geom}(p)$ para un parámetro auxiliar $t$:
		\begin{align*}
			M_{Y_1}(t) &= \mathbb{E}[e^{t Y_1}] = \sum_{k=1}^{\infty} e^{t k} P(Y_1 = k) = \sum_{k=1}^{\infty} e^{t k} q^{k-1} p \\
			&= p e^t \sum_{k=1}^{\infty} e^{t(k-1)} q^{k-1} = p e^t \sum_{j=0}^{\infty} (q e^t)^j, \quad \text{donde } j = k-1.
		\end{align*}
		La sumatoria resultante es una serie geométrica pura de razón $\rho = q e^t$. Para que converga de manera finita, se requiere que $|\rho| < 1 \implies q e^t < 1 \implies t < -\ln(q)$. Bajo esta región de convergencia, sumamos:
		\begin{align*}
			M_{Y_1}(t) = p e^t \left(\dfrac{1}{1 - q e^t}\right) = \dfrac{p e^t}{1 - q e^t}.
		\end{align*}
		\item Por el teorema fundamental de convolución e independencia de FGM, si una variable $X = \sum_{i=1}^r Y_i$ es suma de $r$ variables independientes, su FGM es el producto multiplicativo exacto de las FGM individuales:
		\begin{align*}
			M_X(t) &= \prod_{i=1}^r M_{Y_i}(t) = \prod_{i=1}^r \left(\dfrac{p e^t}{1 - q e^t}\right) = \left(\dfrac{p e^t}{1 - q e^t}\right)^r, \quad \text{para } t < -\ln(q).
		\end{align*}
	\end{enumerate}
\end{solucion}

\begin{solucion}
	\label{sol:3.5.8}
	Para el Problema \ref{prob:3.5.8}:
	\begin{enumerate}
		\item Para obtener la esperanza matemática $\mu_X$, podemos derivar la FGM $M_X(t) = \left(\frac{p e^t}{1 - q e^t}\right)^r$ y evaluarla en $t=0$, o de manera algebraicamente directa por la linealidad del operador esperanza sobre la suma de variables independientes $X = \sum_{i=1}^r Y_i$:
		\begin{align*}
			\mathbb{E}[X] &= \sum_{i=1}^r \mathbb{E}[Y_i] = \sum_{i=1}^r \left(\dfrac{1}{p}\right) = \dfrac{r}{p}.
		\end{align*}
		*(Demostración alternativa por derivación logarítmica de la FGM):* Sea $\ln M_X(t) = r \ln(p) + r t - r \ln(1 - q e^t)$. Derivando respecto a $t$:
		\begin{align*}
			\dfrac{d \ln M_X(t)}{dt} = \dfrac{M_X'(t)}{M_X(t)} = r + \dfrac{r q e^t}{1 - q e^t}.
		\end{align*}
		Evaluando en el origen $t = 0$ donde por definición $M_X(0) = 1$:
		\begin{align*}
			M_X'(0) &= r + \dfrac{r q}{1 - q} = r + \dfrac{r(1-p)}{p} = r + \dfrac{r}{p} - r = \dfrac{r}{p} \implies \mathbb{E}[X] = \dfrac{r}{p}.
		\end{align*}
		\item Por la aditividad de la varianza sobre sumas de variables estocásticamente independientes:
		\begin{align*}
			\text{Var}(X) &= \sum_{i=1}^r \text{Var}(Y_i) = \sum_{i=1}^r \left(\dfrac{q}{p^2}\right) = \dfrac{r q}{p^2} = \dfrac{r(1-p)}{p^2}.
		\end{align*}
		*(Verificación por segunda derivada logarítmica):* Derivando por segunda vez el cumulante $\kappa(t) = \ln M_X(t)$:
		\begin{align*}
			\kappa''(t) &= \dfrac{d}{dt} \left[ r + \dfrac{r q e^t}{1 - q e^t} \right] = \dfrac{r q e^t (1 - q e^t) - (-q e^t) r q e^t}{(1 - q e^t)^2} = \dfrac{r q e^t}{(1 - q e^t)^2}.
		\end{align*}
		Evaluando la segunda derivada del cumulante en el origen $t = 0$, obtenemos de forma unificada y exacta la varianza:
		\begin{align*}
			\text{Var}(X) = \kappa''(0) = \dfrac{r q}{(1 - q)^2} = \dfrac{r q}{p^2}.
		\end{align*}
	\end{enumerate}
\end{solucion}

\begin{solucion}
	\label{sol:3.5.9}
	Para el Problema \ref{prob:3.5.9}:
	\begin{enumerate}
		\item Sea $X$ una variable aleatoria discreta con soporte en $\mathbb{Z}^+ = \set{1, 2, 3, \dots}$ y función de cola de supervivencia de probabilidad definida como $G(k) = P(X > k)$ para $k \ge 0$. Por el axioma de certeza total para el soporte estrictamente positivo, sabemos que $G(0) = P(X > 0) = 1$.
		\item La condición fundamental de pérdida de memoria establece que para todo par de enteros positivos $m, n \in \mathbb{Z}^+$:
		\begin{align*}
			P(X > m+n \mid X > m) = P(X > n) &\implies \dfrac{P(\set{X > m+n} \cap \set{X > m})}{P(X > m)} = P(X > n) \\
			&\implies \dfrac{P(X > m+n)}{P(X > m)} = P(X > n) \\
			&\implies G(m+n) = G(m) G(n).
		\end{align*}
		\item Resolvemos esta ecuación funcional discreta de Cauchy por inducción finita sobre $n$:
		\begin{itemize}
			\item Para $n = 1$: $G(m+1) = G(m) G(1)$ para cualquier entero $m \ge 0$.
			\item Evaluando paso a paso desde el origen:
			\begin{align*}
				G(1) &= G(0+1) = G(0) G(1) = 1 \cdot G(1), \\
				G(2) &= G(1+1) = G(1) G(1) = [G(1)]^2, \\
				G(3) &= G(2+1) = [G(1)]^2 G(1) = [G(1)]^3, \\
				&\vdots \\
				G(k) &= [G(1)]^k, \quad \text{para todo entero } k \ge 0.
			\end{itemize}
		\end{itemize}
		\item Dado que la probabilidad debe decaer para una variable aleatoria propia ($0 < G(1) < 1$), denotamos a la constante de supervivencia de un paso como $q \equiv G(1) = P(X > 1)$. En consecuencia, la función de supervivencia es de forma única $P(X > k) = q^k$.
		\item Finalmente, reconstruimos la función de masa de probabilidad en cada entero puntual $k \ge 1$ mediante la diferencia finita de colas:
		\begin{align*}
			P(X = k) &= P(X > k-1) - P(X > k) = G(k-1) - G(k) \\
			&= q^{k-1} - q^k = q^{k-1}(1 - q).
		\end{align*}
		Definiendo el parámetro de éxito como $p \equiv 1 - q$, obtenemos unívocamente:
		\begin{align*}
			P(X = k) = q^{k-1} p = (1-p)^{k-1} p, \quad \text{para } k \in \set{1, 2, 3, \dots}.
		\end{align*}
		Lo que concluye la demostración formal: en el dominio de las variables aleatorias discretas sobre enteros positivos, la propiedad sin memoria caracteriza de forma exclusiva y necesaria a la **Distribución Geométrica**.
	\end{enumerate}
\end{solucion}

\begin{solucion}
	\label{sol:3.5.10}
	Para el Problema \ref{prob:3.5.10}:
	\begin{enumerate}
		\item Partiendo de la función de masa binomial negativa en su parametrización de conteo de fracasos $Y \sim \text{BN}(r, p)$:
		\begin{align*}
			P(Y = y) = \comb{y+r-1}{y} p^r (1-p)^y, \quad \text{para } y \in \set{0, 1, 2, \dots}.
		\end{align*}
		Bajo el régimen asintótico fijado por $(1-p)r = \lambda$, sustituimos la probabilidad de fracaso $1-p = \frac{\lambda}{r}$ y la de éxito $p = 1 - \frac{\lambda}{r}$:
		\begin{align*}
			P(Y = y) = \comb{y+r-1}{y} \left(1 - \dfrac{\lambda}{r}\right)^r \left(\dfrac{\lambda}{r}\right)^y.
		\end{align*}
		Expandimos el coeficiente combinatorio y reordenamos algebraicamente los factores:
		\begin{align*}
			\comb{y+r-1}{y} \left(\dfrac{\lambda}{r}\right)^y &= \dfrac{(r+y-1)(r+y-2)\cdots(r+1)r}{y!} \cdot \dfrac{\lambda^y}{r^y} \\
			&= \dfrac{\lambda^y}{y!} \left[ \left(\dfrac{r+y-1}{r}\right)\left(\dfrac{r+y-2}{r}\right)\cdots\left(\dfrac{r+1}{r}\right)\left(\dfrac{r}{r}\right) \right].
		\end{align*}
		Observe que dentro del corchete existen exactamente $y$ factores que dependen de $r$. Tomando el límite cuando $r \to \infty$ para un número de fracasos observado $y$ fijo:
		\begin{align*}
			\lim_{r \to \infty} \comb{y+r-1}{y} \left(\dfrac{\lambda}{r}\right)^y &= \dfrac{\lambda^y}{y!} \cdot \left[ 1 \cdot 1 \cdots 1 \cdot 1 \right] = \dfrac{\lambda^y}{y!}.
		\end{align*}
		Por otro lado, por el segundo límite notable de Euler, el factor exponencial acoplado converge a:
		\begin{align*}
			\lim_{r \to \infty} \left(1 - \dfrac{\lambda}{r}\right)^r = e^{-\lambda}.
		\end{align*}
		Multiplicando ambos límites puntuales convergentes, obtenemos de manera exacta la función de probabilidad de Poisson:
		\begin{align*}
			\lim_{r \to \infty} P(Y = y) = \dfrac{\lambda^y e^{-\lambda}}{y!}.
		\end{align*}
		\item **Interpretación Física y Transición de Modelos:**
		La distribución binomial negativa (o ley de Pascal) describe un proceso con una alta tasa de variación individual debido a la espera de un número finito $r$ de eventos macroscópicos. Cuando el número de éxitos exigidos se vuelve infinitamente grande ($r \to \infty$), pero cada ensayo individual es casi invariablemente un éxito ($p \to 1$), la ocurrencia de un fracaso en el transcurso del proceso se convierte en un **evento raro** e independiente.
		
		En este régimen límite, el carácter secuencial y sobredisperso de las sumas geométricas de Pascal ($\sigma^2 > \mu$) se homogeniza perfectamente: la varianza empuja su valor asintótico hacia la media ($\lim_{r \to \infty} \sigma_Y^2 = \lim_{r \to \infty} \frac{\lambda}{p} = \frac{\lambda}{1} = \lambda = \mu_Y$). La ley combinatoria de Pascal colapsa entonces en el **Proceso Puntual de Poisson**, proporcionando el puente teórico unificado entre el conteo discreto de esperas reiteradas y los eventos raros en tiempo continuo.
	\end{enumerate}
\end{solucion}

% ==============================================================================
% SECCIÓN 3.6: DISTRIBUCIÓN HIPERGEOMÉTRICA
% ==============================================================================

\subsubsection*{Nivel 1: Fundamental (Conceptos básicos y directos)}

\begin{problema}
	\label{prob:3.6.1}
	Sea una población finita de tamaño $N$ dividida en $K$ éxitos y $N-K$ fracasos. Se extrae una muestra aleatoria de tamaño $n$ \emph{sin reemplazo}. Sea $X$ el número de éxitos en la muestra.
	\begin{enumerate}
		\item Escriba la función de probabilidad $f_X(k) = P(X = k)$ y determine su soporte exacto $\mathcal{S}_X$.
		\item Demuestre que la suma de probabilidades sobre todo el soporte es igual a $1$ utilizando la identidad de Vandermonde.
	\end{enumerate}
\end{problema}

\begin{sugerencia}
	La identidad de Vandermonde establece que para enteros no negativos $a, b, c$, se cumple $\sum_{j=0}^{c} \comb{a}{j} \comb{b}{c-j} = \comb{a+b}{c}$.
\end{sugerencia}

\begin{problema}
	\label{prob:3.6.2}
	En una planta ensambladora se recibe un lote de $N=30$ tabletas electrónicas. El departamento de control de calidad inspecciona una muestra aleatoria de $n=5$ tabletas sin reemplazo y rechaza todo el lote si encuentra al menos $2$ defectuosas. Si en realidad el lote contiene $K=6$ tabletas defectuosas:
	\begin{enumerate}
		\item Calcule la probabilidad exacta de no encontrar ninguna tableta defectuosa en la muestra ($X = 0$).
		\item Calcule la probabilidad exacta de que el lote sea aceptado ($X \le 1$).
	\end{enumerate}
\end{problema}

\begin{sugerencia}
	El lote se acepta si y solo si $X \in \set{0, 1}$. Utilice la fórmula $P(X=k) = \dfrac{\comb{K}{k}\comb{N-K}{n-k}}{\comb{N}{n}}$.
\end{sugerencia}

\begin{problema}
	\label{prob:3.6.3}
	Sea $X \sim \text{Hiper}(N, K, n)$ el número de éxitos al extraer $n$ elementos sin reemplazo. Podemos expresar a $X$ como la suma de variables indicadoras $X = \sum_{i=1}^n I_i$, donde $I_i = 1$ si la $i$-ésima extracción resulta en éxito y $I_i = 0$ en caso contrario.
	\begin{enumerate}
		\item Determine la probabilidad marginal $P(I_i = 1)$ para cualquier $i \in \set{1, 2, \dots, n}$.
		\item Utilizando la linealidad de la esperanza matemática, demuestre que $\mathbb{E}[X] = n \dfrac{K}{N}$.
	\end{enumerate}
\end{problema}

\begin{sugerencia}
	Por simetría incondicional, la probabilidad marginal de que la $i$-ésima extracción sea un éxito no depende del orden $i$ si no se ha observado la trayectoria previa.
\end{sugerencia}

\subsubsection*{Nivel 2: Operativo (Cálculo e implementación técnica)}

\begin{problema}
	\label{prob:3.6.4}
	Considere una variable aleatoria hipergeométrica $X \sim \text{Hiper}(N, K, n)$ con varianza $\text{Var}(X) = n \dfrac{K}{N} \left(1 - \dfrac{K}{N}\right) \left(\dfrac{N-n}{N-1}\right)$.
	\begin{enumerate}
		\item Identifique y explique el significado del término $\dfrac{N-n}{N-1}$, conocido como Factor de Corrección por Población Finita (FPCF).
		\item Compare el valor numérico de $\text{Var}(X)$ frente al modelo binomial correspondiente ($p = K/N$) cuando $N=50, K=10, n=20$. ¿En qué porcentaje se reduce la varianza real respecto a la aproximación con reemplazo?
	\end{enumerate}
\end{problema}

\begin{sugerencia}
	Evalúe el cociente entre la varianza hipergeométrica y la varianza binomial $np(1-p)$.
\end{sugerencia}

\begin{problema}
	\label{prob:3.6.5}
	Un auditor fiscal revisa un padrón de $N=100$ facturas corporativas, entre las cuales existen $K=15$ con irregularidades tributarias severas. El auditor selecciona una muestra aleatoria de $n=10$ facturas sin reemplazo.
	\begin{enumerate}
		\item Calcule la media $\mu_X$ y la desviación estándar $\sigma_X$ del número de facturas irregulares en la muestra.
		\item Determine la probabilidad exacta de que el auditor detecte al menos $3$ facturas irregulares ($P(X \ge 3)$).
	\end{enumerate}
\end{problema}

\begin{sugerencia}
	Para la cola superior, utilice el complemento $P(X \ge 3) = 1 - [P(X=0) + P(X=1) + P(X=2)]$.
\end{sugerencia}

\begin{problema}
	\label{prob:3.6.6}
	Una base de datos municipal contiene $N=2000$ expedientes de zonificación, de los cuales $K=100$ presentan errores de captura. Se extrae una muestra de $n=20$ expedientes sin reemplazo para control interno.
	\begin{enumerate}
		\item Verifique si se cumple la regla empírica $\dfrac{n}{N} < 0.05$.
		\item Calcule $P(X = 2)$ utilizando la fórmula exacta hipergeométrica y compárela con la aproximación binomial $Y \sim \text{Bin}(20, p=0.05)$. Determine el error absoluto $|P(X=2) - P(Y=2)|$.
	\end{enumerate}
\end{problema}

\begin{sugerencia}
	Calcule $p = K/N = 100/2000 = 0.05$. La aproximación binomial asume independencia muestral.
\end{sugerencia}

\subsubsection*{Nivel 3: Analítico (Demostraciones y modelación matemática)}

\begin{problema}
	\label{prob:3.6.7}
	Sea $X \sim \text{Hiper}(N, K, n)$. Utilizando identidades combinatorias directamente sobre la suma de la función de masa de probabilidad:
	\begin{enumerate}
		\item Demuestre que el primer momento factorial es $\mathbb{E}[X] = n \dfrac{K}{N}$ usando la identidad de absorción $k \comb{K}{k} = K \comb{K-1}{k-1}$.
		\item Demuestre que el segundo momento factorial es $\mathbb{E}[X(X-1)] = n(n-1) \dfrac{K(K-1)}{N(N-1)}$.
	\end{enumerate}
\end{problema}

\begin{sugerencia}
	Para el segundo momento factorial, utilice la identidad de absorción de segundo orden $k(k-1)\comb{K}{k} = K(K-1)\comb{K-2}{k-2}$ y reubique los índices del sumatorio combinatorio.
\end{sugerencia}

\begin{problema}
	\label{prob:3.6.8}
	A partir de la descomposición en variables indicadoras $X = \sum_{i=1}^n I_i$ del Problema \ref{prob:3.6.3}, donde $I_i \sim \text{Bern}(K/N)$:
	\begin{enumerate}
		\item Demuestre que para $i \neq j$, la probabilidad conjunta es $P(I_i = 1, I_j = 1) = \dfrac{K(K-1)}{N(N-1)}$.
		\item Deduzca que la covarianza entre dos extracciones distintas sin reemplazo es negativa y está dada por $\text{Cov}(I_i, I_j) = -\dfrac{K(N-K)}{N^2(N-1)}$.
		\item Utilizando la fórmula general de la varianza para sumas de variables correlacionadas, obtenga la expresión exacta de $\text{Var}(X)$.
	\end{enumerate}
\end{problema}

\begin{sugerencia}
	$\text{Var}\left(\sum_{i=1}^n I_i\right) = \sum_{i=1}^n \text{Var}(I_i) + \sum_{i \neq j} \text{Cov}(I_i, I_j)$. Exprese el resultado en función del FPCF $\frac{N-n}{N-1}$.
\end{sugerencia}

\subsubsection*{Nivel 4: Desafiante (Problemas avanzados o retos de ingeniería)}

\begin{problema}
	\label{prob:3.6.9}
	Demuestre formalmente el límite combinatorio asintótico que conecta el muestreo sin reemplazo con el muestreo con reemplazo:
	Sea $X_N \sim \text{Hiper}(N, K_N, n)$ de modo que cuando $N \to \infty$ y $K_N \to \infty$, la proporción de éxitos en la población tiende a una constante fija $p = \lim_{N \to \infty} \dfrac{K_N}{N} \in (0, 1)$, mientras que el tamaño de muestra $n$ permanece constante. Demuestre puntualmente que para todo $k \in \set{0, 1, \dots, n}$:
	\begin{align*}
		\lim_{N \to \infty} P(X_N = k) = \comb{n}{k} p^k (1-p)^{n-k}.
	\end{align*}
\end{problema}

\begin{sugerencia}
	Descomponga los coeficientes combinatorios en factoriales decrecientes: $\comb{K}{k}\comb{N-K}{n-k}/\comb{N}{n} = \comb{n}{k} \frac{K^{(k)} (N-K)^{(n-k)}}{N^{(n)}}$ y tome el límite de cada cociente lineal.
\end{sugerencia}

\begin{problema}
	\label{prob:3.6.10}
	**Prueba Exacta de Fisher en Ensayos Clínicos:**
	En un estudio clínico sobre un nuevo agente antitrombótico, se seleccionaron $N=25$ pacientes en total, de los cuales $K=12$ recibieron la terapia activa y $N-K=13$ recibieron placebo. Al finalizar el estudio, se observó que exactamente $n=8$ pacientes experimentaron la formación de trombos en general, mientras que $17$ no presentaron trombos.
	
	Bajo la hipótesis nula ($H_0$) de que el tratamiento no tiene ningún efecto y la trombosis ocurre independientemente del grupo al que fue asignado el paciente:
	\begin{enumerate}
		\item Demuestre teóricamente que, condicionado a los marginales fijos de la tabla de contingencia $2 \times 2$, la variable aleatoria $X$ (número de pacientes con trombosis en el grupo de terapia activa) sigue una distribución hipergeométrica con parámetros $N=25, K=12, n=8$.
		\item Si en el estudio se observó que solo $X = 1$ paciente del grupo de terapia activa desarrolló trombos (frente a $7$ en el placebo), calcule el *valor p de una cola inferior* ($P(X \le 1)$) e interprete estadísticamente si existe evidencia de eficacia terapéutica al nivel $\alpha = 0.05$.
	\end{enumerate}
\end{problema}

\begin{sugerencia}
	El valor $p$ exacto de Fisher de una cola en la dirección observada es la suma hipergeométrica $\sum_{j=0}^{1} P(X=j)$.
\end{sugerencia}

% --- SOLUCIONES DE LA SECCIÓN 3.6 ---

\begin{solucion}
	\label{sol:3.6.1}
	Para el Problema \ref{prob:3.6.1}:
	\begin{enumerate}
		\item El número de formas de elegir $k$ éxitos de un total de $K$ disponibles es $\comb{K}{k}$, y por cada una de estas formas, se deben elegir $(n-k)$ fracasos de los $(N-K)$ disponibles, con $\comb{N-K}{n-k}$ posibilidades. El número total de configuraciones equiprobables al extraer $n$ elementos de $N$ sin orden es $\comb{N}{n}$. Así:
		\begin{align*}
			f_X(k) = P(X=k) = \dfrac{\comb{K}{k}\comb{N-K}{n-k}}{\comb{N}{n}}.
		\end{align*}
		El soporte $\mathcal{S}_X$ requiere que los argumentos superiores sean mayores o iguales que los inferiores en cada coeficiente: $k \le K$ y $n-k \le N-K \implies k \ge n-(N-K)$. Además, $k \ge 0$ y $k \le n$. Por tanto, el intervalo exacto es:
		\begin{align*}
			\mathcal{S}_X = \set{\max(0, n-(N-K)), \dots, \min(n, K)}.
		\end{align*}
		\item Sumando la función de probabilidad sobre todas las configuraciones posibles y aplicando la identidad combinatoria de Vandermonde en el numerador:
		\begin{align*}
			\sum_{k \in \mathcal{S}_X} f_X(k) &= \dfrac{1}{\comb{N}{n}} \sum_{k \in \mathcal{S}_X} \comb{K}{k}\comb{N-K}{n-k} \\
			&= \dfrac{1}{\comb{N}{n}} \comb{K + (N-K)}{k + (n-k)} = \dfrac{\comb{N}{n}}{\comb{N}{n}} = 1.
		\end{align*}
	\end{enumerate}
\end{solucion}

\begin{solucion}
	\label{sol:3.6.2}
	Para el Problema \ref{prob:3.6.2}:
	Aquí los parámetros son $N=30$, $K=6$ (defectuosas), $N-K=24$ (buenas) y $n=5$ (muestra).
	\begin{enumerate}
		\item La probabilidad de que las $5$ tabletas muestreadas sean buenas ($k=0$ defectuosas):
		\begin{align*}
			P(X = 0) &= \dfrac{\comb{6}{0}\comb{24}{5}}{\comb{30}{5}} = \dfrac{1 \times 42504}{142506} \approx 0.298259 \approx 0.2983.
		\end{align*}
		\item La probabilidad de encontrar exactamente $k=1$ tableta defectuosa:
		\begin{align*}
			P(X = 1) &= \dfrac{\comb{6}{1}\comb{24}{4}}{\comb{30}{5}} = \dfrac{6 \times 10626}{142506} = \dfrac{63756}{142506} \approx 0.447391 \approx 0.4474.
		\end{align*}
		La probabilidad exacta de que el lote sea aceptado ($X \le 1$) es la suma puntual:
		\begin{align*}
			P(\text{Aceptación}) &= P(X = 0) + P(X = 1) \approx 0.298259 + 0.447391 = 0.74565 \approx 0.7457.
		\end{align*}
		Existe un $74.57\%$ de probabilidad de aceptar este lote con 6 piezas defectuosas. La probabilidad de rechazarlo es $1 - 0.7457 = 25.43\%$.
	\end{enumerate}
\end{solucion}

\begin{solucion}
	\label{sol:3.6.3}
	Para el Problema \ref{prob:3.6.3}:
	\begin{enumerate}
		\item La variable indicadora $I_i$ vale $1$ si la $i$-ésima tableta/elemento extraído es un éxito. Por simetría incondicional (antes de conocer los resultados de las extracciones previas o posteriores), cada uno de los $N$ elementos de la población tiene la misma probabilidad de ocupar la $i$-ésima posición en la muestra:
		\begin{align*}
			P(I_i = 1) = \dfrac{K}{N}, \quad \forall i \in \set{1, 2, \dots, n}.
		\end{align*}
		\item La esperanza de una variable indicadora es su propia probabilidad de éxito: $\mathbb{E}[I_i] = P(I_i = 1) = \dfrac{K}{N}$. Por el teorema de linealidad de la esperanza (el cual es válido incluso para variables aleatorias dependientes o correlacionadas):
		\begin{align*}
			\mathbb{E}[X] &= \mathbb{E}\left[\sum_{i=1}^n I_i\right] = \sum_{i=1}^n \mathbb{E}[I_i] = \sum_{i=1}^n \dfrac{K}{N} = n \dfrac{K}{N}.
		\end{align*}
	\end{enumerate}
\end{solucion}

\begin{solucion}
	\label{sol:3.6.4}
	Para el Problema \ref{prob:3.6.4}:
	\begin{enumerate}
		\item El término $\text{FPCF} = \dfrac{N-n}{N-1}$ mide la reducción geométrica en la variabilidad muestral debida a que cada extracción sin reemplazo disminuye el conjunto disponible. 
		- Cuando $n=1$, el factor vale $\frac{N-1}{N-1} = 1$ (coincidiendo con Bernoulli).
		- Cuando $n=N$ (censo total), el factor vale $\frac{0}{N-1} = 0$, ya que al inspeccionar toda la población no existe ninguna incertidumbre muestral ($\text{Var}(X) = 0$).
		\item Con los parámetros $N=50, K=10, n=20$:
		- Varianza Binomial (con reemplazo, $p = 10/50 = 0.20$):
		  \begin{align*}
		  	\text{Var}_{\text{Bin}}(X) = np(1-p) = 20(0.20)(0.80) = 3.20.
		  \end{align*}
		- Factor de Corrección por Población Finita:
		  \begin{align*}
		  	\text{FPCF} = \dfrac{50-20}{50-1} = \dfrac{30}{49} \approx 0.612245.
		  \end{align*}
		- Varianza Hipergeométrica real:
		  \begin{align*}
		  	\text{Var}_{\text{Hiper}}(X) = 3.20 \times \dfrac{30}{49} = \dfrac{96}{49} \approx 1.9592.
		  \end{align*}
		La varianza se reduce en un $1 - 0.612245 = 38.78\%$ gracias a la información no redundante capturada por el muestreo sin reemplazo.
	\end{enumerate}
\end{solucion}

\begin{solucion}
	\label{sol:3.6.5}
	Para el Problema \ref{prob:3.6.5}:
	Tenemos $N=100$, $K=15$ y $n=10$.
	\begin{enumerate}
		\item Cálculo de momentos de primer y segundo orden:
		\begin{align*}
			\mu_X &= 10 \left(\dfrac{15}{100}\right) = 1.50. \\
			\sigma_X^2 &= 10(0.15)(0.85)\left(\dfrac{100-10}{100-1}\right) = 1.275 \left(\dfrac{90}{99}\right) = 1.275 \times \dfrac{10}{11} \approx 1.15909. \\
			\sigma_X &= \sqrt{1.15909} \approx 1.0766.
		\end{align*}
		\item Para hallar $P(X \ge 3)$, calculamos los tres primeros términos del soporte ($k=0, 1, 2$):
		\begin{align*}
			P(X=0) &= \dfrac{\comb{15}{0}\comb{85}{10}}{\comb{100}{10}} \approx 0.18684, \\
			P(X=1) &= \dfrac{\comb{15}{1}\comb{85}{9}}{\comb{100}{10}} \approx 0.36402, \\
			P(X=2) &= \dfrac{\comb{15}{2}\comb{85}{8}}{\comb{100}{10}} \approx 0.28634.
		\end{align*}
		Sumando las probabilidades inferiores: $P(X \le 2) \approx 0.18684 + 0.36402 + 0.28634 = 0.83720$.
		Por la regla del complemento:
		\begin{align*}
			P(X \ge 3) = 1 - P(X \le 2) \approx 1 - 0.83720 = 0.1628.
		\end{align*}
		El auditor tiene un $16.28\%$ de probabilidad de detectar $3$ o más facturas irregulares en una muestra de solo $10$ documentos.
	\end{enumerate}
\end{solucion}

\begin{solucion}
	\label{sol:3.6.6}
	Para el Problema \ref{prob:3.6.6}:
	\begin{enumerate}
		\item El cociente muestral es $\dfrac{n}{N} = \dfrac{20}{2000} = 0.01 = 1\%$. Como $1\% < 5\%$, se cumple ampliamente el criterio práctico de independencia aproximada.
		\item Probabilidad exacta Hipergeométrica ($N=2000, K=100, n=20, k=2$):
		\begin{align*}
			P(X=2) = \dfrac{\comb{100}{2}\comb{1900}{18}}{\comb{2000}{20}} \approx 0.189725.
		\end{align*}
		Probabilidad Binomial aproximada ($n=20, p=100/2000=0.05, k=2$):
		\begin{align*}
			P(Y=2) = \comb{20}{2}(0.05)^2(0.95)^{18} = 190(0.0025)(0.397214) \approx 0.188677.
		\end{align*}
		El error absoluto exacto entre ambos modelos es:
		\begin{align*}
			|P(X=2) - P(Y=2)| \approx |0.189725 - 0.188677| = 0.001048 \approx 0.105\%.
		\end{align*}
		La aproximación binomial introduce un error marginal de apenas una décima de punto porcentual.
	\end{enumerate}
\end{solucion}

\begin{solucion}
	\label{sol:3.6.7}
	Para el Problema \ref{prob:3.6.7}:
	\begin{enumerate}
		\item Aplicando la identidad combinatoria $k \comb{K}{k} = K \comb{K-1}{k-1}$ sobre la definición de esperanza:
		\begin{align*}
			\mathbb{E}[X] &= \sum_{k=1}^n k \dfrac{\comb{K}{k}\comb{N-K}{n-k}}{\comb{N}{n}} = \dfrac{1}{\comb{N}{n}} \sum_{k=1}^n K \comb{K-1}{k-1}\comb{N-K}{n-k}.
		\end{align*}
		Extraemos $K$ y reescribimos el denominador utilizando $n \comb{N}{n} = N \comb{N-1}{n-1} \implies \comb{N}{n} = \frac{N}{n} \comb{N-1}{n-1}$:
		\begin{align*}
			\mathbb{E}[X] &= n \dfrac{K}{N} \sum_{k=1}^n \dfrac{\comb{K-1}{k-1}\comb{N-K}{n-k}}{\comb{N-1}{n-1}}.
		\end{align*}
		Con el cambio de variable $j = k-1$, la suma sobre $j \in \set{0, \dots, n-1}$ de $\comb{K-1}{j}\comb{(N-1)-(K-1)}{(n-1)-j}/\comb{N-1}{n-1}$ es la suma de una PMF hipergeométrica con parámetros $(N-1, K-1, n-1)$, cuyo total es exactamente $1$. Queda demostrado $\mathbb{E}[X] = n \dfrac{K}{N}$.
		\item Para el segundo momento factorial, aplicamos $k(k-1)\comb{K}{k} = K(K-1)\comb{K-2}{k-2}$:
		\begin{align*}
			\mathbb{E}[X(X-1)] &= \sum_{k=2}^n k(k-1) \dfrac{\comb{K}{k}\comb{N-K}{n-k}}{\comb{N}{n}} \\
			&= \dfrac{K(K-1)}{\comb{N}{n}} \sum_{k=2}^n \comb{K-2}{k-2}\comb{N-K}{n-k}.
		\end{align*}
		Usando $\comb{N}{n} = \dfrac{N(N-1)}{n(n-1)}\comb{N-2}{n-2}$:
		\begin{align*}
			\mathbb{E}[X(X-1)] &= n(n-1)\dfrac{K(K-1)}{N(N-1)} \sum_{j=0}^{n-2} \dfrac{\comb{K-2}{j}\comb{(N-2)-(K-2)}{(n-2)-j}}{\comb{N-2}{n-2}} \\
			&= n(n-1)\dfrac{K(K-1)}{N(N-1)} \times 1 = n(n-1)\dfrac{K(K-1)}{N(N-1)}.
		\end{align*}
	\end{enumerate}
\end{solucion}

\begin{solucion}
	\label{sol:3.6.8}
	Para el Problema \ref{prob:3.6.8}:
	\begin{enumerate}
		\item Para dos extracciones distintas $i \neq j$, el evento $(I_i = 1, I_j = 1)$ significa que la $i$-ésima y la $j$-ésima extracción son éxitos. Por regla de la multiplicación incondicional:
		\begin{align*}
			P(I_i = 1, I_j = 1) &= P(I_i = 1) P(I_j = 1 \mid I_i = 1) = \dfrac{K}{N} \left(\dfrac{K-1}{N-1}\right).
		\end{align*}
		\item La covarianza entre dos variables indicadoras se obtiene por definición:
		\begin{align*}
			\text{Cov}(I_i, I_j) &= \mathbb{E}[I_i I_j] - \mathbb{E}[I_i]\mathbb{E}[I_j] = P(I_i = 1, I_j = 1) - \left(\dfrac{K}{N}\right)^2 \\
			&= \dfrac{K(K-1)}{N(N-1)} - \dfrac{K^2}{N^2} = \dfrac{K}{N} \left[ \dfrac{K-1}{N-1} - \dfrac{K}{N} \right] \\
			&= \dfrac{K}{N} \left[ \dfrac{N(K-1) - K(N-1)}{N(N-1)} \right] = \dfrac{K}{N} \left[ \dfrac{NK - N - NK + K}{N(N-1)} \right] \\
			&= -\dfrac{K(N-K)}{N^2(N-1)} < 0.
		\end{align*}
		\item Varianza de la suma de indicadoras correlacionadas:
		\begin{align*}
			\text{Var}(X) &= \sum_{i=1}^n \text{Var}(I_i) + \sum_{i \neq j} \text{Cov}(I_i, I_j).
		\end{align*}
		Como $\text{Var}(I_i) = \frac{K}{N}\left(1 - \frac{K}{N}\right)$ y existen $n(n-1)$ pares ordenados $i \neq j$ con la misma covarianza:
		\begin{align*}
			\text{Var}(X) &= n \dfrac{K}{N}\left(\dfrac{N-K}{N}\right) - n(n-1) \dfrac{K(N-K)}{N^2(N-1)} \\
			&= n \dfrac{K}{N}\left(\dfrac{N-K}{N}\right) \left[ 1 - \dfrac{n-1}{N-1} \right] \\
			&= n \dfrac{K}{N}\left(1 - \dfrac{K}{N}\right) \left[ \dfrac{(N-1) - (n-1)}{N-1} \right] = n \dfrac{K}{N}\left(1 - \dfrac{K}{N}\right) \left(\dfrac{N-n}{N-1}\right).
		\end{align*}
	\end{enumerate}
\end{solucion}

\begin{solucion}
	\label{sol:3.6.9}
	Para el Problema \ref{prob:3.6.9}:
	Expresamos la función de masa hipergeométrica utilizando factoriales decrecientes $A^{(m)} = A(A-1)\cdots(A-m+1)$:
	\begin{align*}
		P(X_N = k) &= \dfrac{\comb{K_N}{k}\comb{N-K_N}{n-k}}{\comb{N}{n}} = \dfrac{\dfrac{K_N^{(k)}}{k!} \cdot \dfrac{(N-K_N)^{(n-k)}}{(n-k)!}}{\dfrac{N^{(n)}}{n!}} = \comb{n}{k} \dfrac{K_N^{(k)} (N-K_N)^{(n-k)}}{N^{(n)}}.
	\end{align*}
	Para el factor con $K_N^{(k)}$, dividimos cada uno de los $k$ factores decrecientes entre $N$:
	\begin{align*}
		\dfrac{K_N^{(k)}}{N^k} &= \left(\dfrac{K_N}{N}\right)\left(\dfrac{K_N-1}{N}\right)\cdots\left(\dfrac{K_N-k+1}{N}\right) \xrightarrow[N \to \infty]{} p \cdot p \cdots p = p^k.
	\end{align*}
	De manera idéntica, para los $(n-k)$ factores del segundo término, como $\frac{N-K_N}{N} = 1 - \frac{K_N}{N} \to 1-p$:
	\begin{align*}
		\dfrac{(N-K_N)^{(n-k)}}{N^{n-k}} &\xrightarrow[N \to \infty]{} (1-p)^{n-k}.
	\end{align*}
	Finalmente, el denominador $N^{(n)}/N^n \to 1$ ya que $n$ es una constante finita:
	\begin{align*}
		\lim_{N \to \infty} P(X_N = k) &= \comb{n}{k} \left[ \lim_{N \to \infty} \dfrac{K_N^{(k)}}{N^k} \right] \left[ \lim_{N \to \infty} \dfrac{(N-K_N)^{(n-k)}}{N^{n-k}} \right] \left[ \lim_{N \to \infty} \dfrac{N^n}{N^{(n)}} \right] \\
		&= \comb{n}{k} p^k (1-p)^{n-k}.
	\end{align*}
\end{solucion}

\begin{solucion}
	\label{sol:3.6.10}
	Para el Problema \ref{prob:3.6.10}:
	\begin{enumerate}
		\item En la tabla $2 \times 2$, tenemos $N=25$ sujetos con totales marginales de tratamientos fijos ($K=12$ en terapia activa, $N-K=13$ en placebo) y totales marginales de respuesta fijos ($n=8$ con trombos, $N-n=17$ sin trombos).
		
		Bajo la hipótesis nula $H_0$ de independencia (los $n=8$ eventos trombóticos se distribuyen al azar entre los $25$ pacientes), cada combinación de elegir los $8$ pacientes trombóticos entre los $N$ disponibles tiene probabilidad $1/\comb{N}{n}$. El número de formas en que exactamente $k$ de esos $8$ trombos recaigan en el grupo activo (de tamaño $K$) y los restantes $(n-k)$ recaigan en el grupo placebo (de tamaño $N-K$) es $\comb{K}{k}\comb{N-K}{n-k}$. Por tanto:
		\begin{align*}
			P(X = k \mid \text{marginales fijos}) = \dfrac{\comb{K}{k}\comb{N-K}{n-k}}{\comb{N}{n}} \sim \text{Hiper}(N=25, K=12, n=8).
		\end{align*}
		\item Evaluamos la probabilidad exacta del resultado observado ($X=1$) y los más extremos en favor de la eficacia terapéutica ($X=0$):
		\begin{align*}
			P(X = 0) &= \dfrac{\comb{12}{0}\comb{13}{8}}{\comb{25}{8}} = \dfrac{1 \times 1287}{1081575} \approx 0.001190. \\
			P(X = 1) &= \dfrac{\comb{12}{1}\comb{13}{7}}{\comb{25}{8}} = \dfrac{12 \times 1716}{1081575} = \dfrac{20592}{1081575} \approx 0.019039.
		\end{align*}
		El *valor p de una cola inferior* de la Prueba Exacta de Fisher es:
		\begin{align*}
			p\text{-value} = P(X \le 1) = P(X=0) + P(X=1) \approx 0.001190 + 0.019039 = 0.020229 \approx 0.0202.
		\end{align*}
		Como $p\text{-value} \approx 0.0202 < \alpha = 0.05$, se rechaza la hipótesis nula $H_0$. Existe evidencia estadística estadísticamente significativa de que la terapia activa reduce la incidencia de trombosis respecto al placebo.
	\end{enumerate}
\end{solucion}

% ==============================================================================
% SECCIÓN 03.07: DISTRIBUCIÓN DE POISSON Y PROCESOS DE LLEGADA
% ==============================================================================

\subsection*{Enunciados de los problemas --- Sección 03.07}

\subsubsection*{Nivel Fundamental}

\begin{problema}
	\label{prob:3.7.1}
	Un centro de atención telefónica recibe llamadas de manera independiente a una tasa promedio constante de $\lambda = 3$ llamadas por minuto. Sea $X$ la variable aleatoria que cuenta el número de llamadas recibidas en un intervalo fijo de un minuto.
	\begin{enumerate}
		\item Especifique formalmente la función de masa de probabilidad $P(X = k)$ utilizando la fórmula de Poisson $P(X = k) = \dfrac{e^{-\lambda}\lambda^{k}}{k!}$ para $k \in \set{0, 1, 2, \dots}$.
		\item Calcule explícitamente las probabilidades puntuales $P(X = 0)$, $P(X = 1)$ y $P(X = 2)$ y la probabilidad acumulada $P(X \ge 4)$.
	\end{enumerate}
\end{problema}

\begin{problema}
	\label{prob:3.7.2}
	En una planta de manufactura de componentes electrónicos, el número de defectos por lote sigue una distribución de Poisson con parámetro $\lambda = 2.5$ defectos por lote.
	\begin{enumerate}
		\item Calcule la esperanza matemática $\E[X]$ y la varianza $\Var(X)$ mediante la suma directa sobre el soporte $k \in \set{0, 1, 2, \dots, 15}$.
		\item Verifique numéricamente la propiedad de equidispersión $\E[X] = \Var(X) = \lambda$ y comente su importancia práctica en el control estadístico de procesos.
	\end{enumerate}
\end{problema}

\begin{problema}
	\label{prob:3.7.3}
	El número de personas que llegan por día a la sala de urgencias de un hospital sigue una distribución de Poisson con media $\lambda = 5$ personas por día.
	\begin{enumerate}
		\item Calcule la probabilidad de que a lo sumo lleguen $3$ personas en un día dado, es decir, $P(X \le 3) = \sum_{k=0}^{3} \dfrac{5^{k} e^{-5}}{k!}$.
		\item Calcule la probabilidad de que lleguen al menos $8$ personas en un día, $P(X \ge 8) = 1 - P(X \le 7)$.
	\end{enumerate}
\end{problema}

\subsubsection*{Nivel Operativo}

\begin{problema}
	\label{prob:3.7.4}
	En un ensayo clínico de fase III se monitorean $n = 2000$ pacientes durante un año para detectar una reacción adversa extremadamente rara con probabilidad verdadera $p = 0.001$ por paciente-año (de modo que $\lambda = np = 2$).
	\begin{enumerate}
		\item Calcule la probabilidad exacta de observar exactamente $k = 2$ reacciones adversas bajo el modelo binomial $Y \sim \text{Bin}(2000, 0.001)$.
		\item Calcule la aproximación de Poisson $P(X = 2) = \dfrac{e^{-2} 2^{2}}{2!}$ y determine el error absoluto $\abs{P_{\text{Bin}}(Y=2) - P_{\text{Pois}}(X=2)}$.
		\item Justifique teóricamente por qué la aproximación de Poisson es excelente bajo la regla operativa $n \ge 50$ y $p \le 0.05$ (con $np \le 5$).
	\end{enumerate}
\end{problema}

\begin{problema}
	\label{prob:3.7.5}
	Un servidor web recibe en promedio $\lambda = 8$ solicitudes HTTP por segundo, modeladas como un proceso de Poisson.
	\begin{enumerate}
		\item Calcule la probabilidad de que en un segundo particular el servidor no reciba ninguna solicitud.
		\item Calcule la probabilidad de que el servidor sufra saturación recibiendo $\ge 15$ solicitudes en un segundo.
		\item Halle el cuantil $q_{0.95}$, es decir, el menor entero $q$ tal que $P(X \le q) \ge 0.95$.
	\end{enumerate}
\end{problema}

\begin{problema}
	\label{prob:3.7.6}
	Un sistema de monitoreo industrial registra fallas en una máquina crítica con una tasa promedio de $\lambda = 0.5$ fallas por hora, bajo un proceso de Poisson.
	\begin{enumerate}
		\item Calcule la probabilidad de que ocurra al menos una falla en un turno de $4$ horas consecutivas.
		\item Compare el resultado anterior con el que se obtendría al modelar el número de fallas en $4$ horas como una binomial con $n = 4$ ensayos horarios de Bernoulli con $p$ equivalente.
	\end{enumerate}
\end{problema}

\subsubsection*{Nivel Analítico}

\begin{problema}
	\label{prob:3.7.7}
	Sea $X \sim \text{Pois}(\lambda)$ una variable aleatoria con distribución de Poisson y parámetro $\lambda > 0$.
	\begin{enumerate}
		\item Demuestre analíticamente que $\E[X] = \lambda$ partiendo de la identidad de la serie exponencial $e^{\lambda} = \sum_{k=0}^{\infty} \dfrac{\lambda^{k}}{k!}$ y derivando término a término.
		\item Demuestre analíticamente que $\E[X^{2}] = \lambda^{2} + \lambda$ utilizando la derivada segunda de la misma serie exponencial.
		\item Concluya que $\Var(X) = \E[X^{2}] - (\E[X])^{2} = \lambda$, confirmando la propiedad de equidispersión.
	\end{enumerate}
\end{problema}

\begin{problema}
	\label{prob:3.7.8}
	Considere dos variables aleatorias independientes $X \sim \text{Pois}(\lambda_{1})$ y $Y \sim \text{Pois}(\lambda_{2})$.
	\begin{enumerate}
		\item Demuestre rigurosamente, mediante convolución discreta, que la suma $Z = X + Y$ sigue una distribución $\text{Pois}(\lambda_{1} + \lambda_{2})$. Sugerencia: aplique el Teorema del Binomio de Newton a $(e^{-\lambda_{1}} e^{-\lambda_{2}})$.
		\item Verifique que la propiedad de aditividad es coherente con la conservación de la esperanza $\E[Z] = \lambda_{1} + \lambda_{2}$ y la varianza $\Var(Z) = \lambda_{1} + \lambda_{2}$.
	\end{enumerate}
\end{problema}

\subsubsection*{Nivel Desafiante}

\begin{problema}
	\label{prob:3.7.9}
	Sean $X \sim \text{Pois}(\lambda_{1})$ y $Y \sim \text{Pois}(\lambda_{2})$ dos variables aleatorias independientes que representan el número de eventos atendidos por dos servidores web en un mismo intervalo de tiempo.
	\begin{enumerate}
		\item Demuestre analíticamente que la distribución condicional de $X$ dado el evento $(X + Y = n)$ es una distribución binomial con parámetros $n$ y $p = \lambda_{1} / (\lambda_{1} + \lambda_{2})$:
		\begin{align*}
			P(X = k \mid X + Y = n) = \dfrac{P(X = k) P(Y = n - k)}{P(X + Y = n)} = \comb{n}{k} \left(\dfrac{\lambda_{1}}{\lambda_{1} + \lambda_{2}}\right)^{k} \left(\dfrac{\lambda_{2}}{\lambda_{1} + \lambda_{2}}\right)^{n-k}.
		\end{align*}
		\item Aplique el resultado a un caso donde $\lambda_{1} = 3$ y $\lambda_{2} = 5$. Si en una hora se atendieron en total $n = 8$ solicitudes entre ambos servidores, calcule la probabilidad de que el servidor $1$ haya atendido exactamente $k = 3$ de ellas.
	\end{enumerate}
\end{problema}

\begin{problema}
	\label{prob:3.7.10}
	En un proceso espacial de Poisson, una lámina metálica de $A = 10\,\text{m}^{2}$ presenta defectos distribuidos uniformemente con una densidad $\lambda = 0.5$ defectos por $\text{m}^{2}$.
	\begin{enumerate}
		\item Modele el número total de defectos $X$ en la lámina completa y calcule $P(X = 0)$ y $P(X \ge 8)$.
		\item La lámina se divide en dos regiones disjuntas $A_{1} = 4\,\text{m}^{2}$ y $A_{2} = 6\,\text{m}^{2}$. Sean $X_{1}$ y $X_{2}$ los defectos en cada región. Demuestre que $X_{1} \sim \text{Pois}(2)$ y $X_{2} \sim \text{Pois}(3)$ son independientes.
		\item Demuestre que, condicionalmente al total $X_{1} + X_{2} = n$, la distribución de $X_{1}$ es $\text{Bin}(n, 4/10)$, resultado análogo al del Problema \ref{prob:3.7.9} pero con $p = A_{1}/A$.
	\end{enumerate}
\end{problema}

\subsection*{Sugerencias de los problemas --- Sección 03.07}

\begin{sugerencia}
	\label{sug:3.7.1}
	Para el Problema \ref{prob:3.7.1}, evalúe la fórmula $P(X=k) = e^{-\lambda}\lambda^{k}/k!$ directamente con $\lambda = 3$ y $k \in \set{0, 1, 2}$. Para $P(X \ge 4)$, use el complemento $1 - \sum_{k=0}^{3} P(X=k)$.
\end{sugerencia}

\begin{sugerencia}
	\label{sug:3.7.2}
	Para el Problema \ref{prob:3.7.2}, evalúe las sumas $\E[X] = \sum_{k=0}^{15} k \cdot P(X=k)$ y $\E[X^{2}] = \sum_{k=0}^{15} k^{2} \cdot P(X=k)$. Verifique que la diferencia $\E[X^{2}] - (\E[X])^{2}$ se aproxima a $2.5$ con tolerancia de $10^{-3}$.
\end{sugerencia}

\begin{sugerencia}
	\label{sug:3.7.3}
	Para el Problema \ref{prob:3.7.3}, use la serie truncada $P(X \le 3) = e^{-5}\left(\dfrac{5^{0}}{0!} + \dfrac{5^{1}}{1!} + \dfrac{5^{2}}{2!} + \dfrac{5^{3}}{3!}\right)$ y análogamente para $P(X \le 7)$. El factor común $e^{-5} \approx 0.006738$.
\end{sugerencia}

\begin{sugerencia}
	\label{sug:3.7.4}
	Para el Problema \ref{prob:3.7.4}, use `scipy.stats.binom.pmf(2, 2000, 0.001)` y `scipy.stats.poisson.pmf(2, mu=2)`. La diferencia debería ser menor que $5 \times 10^{-4}$ gracias a la regla $n \cdot p^{2} = 2000 \cdot 10^{-6} = 0.002 \ll 1$.
\end{sugerencia}

\begin{sugerencia}
	\label{sug:3.7.5}
	Para el Problema \ref{prob:3.7.5}, use el CDF de Poisson `stats.poisson.cdf(q, mu=8)` evaluado en distintos $q$ para encontrar el cuantil. Itere desde $q = 10$ hasta que el CDF supere $0.95$.
\end{sugerencia}

\begin{sugerencia}
	\label{sug:3.7.6}
	Para el Problema \ref{prob:3.7.6}, el número de fallas en $4$ horas sigue una Poisson con $\lambda' = 4 \cdot 0.5 = 2$ (por escalamiento lineal del tiempo). Calcule $P(X \ge 1) = 1 - e^{-2}$. En el modelo binomial equivalente, use $n=4$ ensayos horarios y $p = 1 - e^{-0.5} \approx 0.3935$.
\end{sugerencia}

\begin{sugerencia}
	\label{sug:3.7.7}
	Para el Problema \ref{prob:3.7.7}, multiplique ambos lados de la identidad $e^{\lambda} = \sum_{k=0}^{\infty} \lambda^{k}/k!$ por $e^{-\lambda}$ para obtener la normalización. Derive una vez respecto a $\lambda$ y multiplique por $\lambda$ para obtener $\E[X]$. Derive dos veces para $\E[X(X-1)] = \lambda^{2}$ y use $\E[X^{2}] = \E[X(X-1)] + \E[X]$.
\end{sugerencia}

\begin{sugerencia}
	\label{sug:3.7.8}
	Para el Problema \ref{prob:3.7.8}, exprese la convolución $P(Z = m) = \sum_{k=0}^{m} P(X=k) P(Y=m-k) = \sum_{k=0}^{m} \dfrac{e^{-\lambda_{1}} \lambda_{1}^{k}}{k!} \cdot \dfrac{e^{-\lambda_{2}} \lambda_{2}^{m-k}}{(m-k)!}$. Factorice $e^{-(\lambda_{1}+\lambda_{2})}$ y reconozca el Teorema del Binomio de Newton: $\sum_{k=0}^{m} \comb{m}{k} \lambda_{1}^{k} \lambda_{2}^{m-k} = (\lambda_{1} + \lambda_{2})^{m}$.
\end{sugerencia}

\begin{sugerencia}
	\label{sug:3.7.9}
	Para el Problema \ref{prob:3.7.9}, sustituya las PMFs de Poisson en el numerador y denominador, factorice $e^{-(\lambda_{1}+\lambda_{2})}(\lambda_{1}+\lambda_{2})^{n}/n!$ en el denominador (por el Problema \ref{prob:3.7.8}), y observe que el cociente colapsa a una PMF binomial con $p = \lambda_{1}/(\lambda_{1}+\lambda_{2})$.
\end{sugerencia}

\begin{sugerencia}
	\label{sug:3.7.10}
	Para el Problema \ref{prob:3.7.10}, use que en un proceso espacial de Poisson con intensidad $\lambda$ por unidad de área, una región de área $A$ contiene $X \sim \text{Pois}(\lambda A)$ defectos. Para (b), las regiones disjuntas son independientes. Para (c), use el resultado del Problema \ref{prob:3.7.9} con $p = A_{1}/(A_{1}+A_{2}) = 4/10 = 0.4$.
\end{sugerencia}

\subsection*{Soluciones de los problemas --- Sección 03.07}

\begin{solucion}
	\label{sol:3.7.1}
	Para el Problema \ref{prob:3.7.1}:
	\begin{enumerate}
		\item El soporte de $X$ es $\mathcal{S}_{X} = \set{0, 1, 2, \dots}$ (enteros no negativos). La PMF canónica es:
		\begin{align*}
			P(X = k) = \dfrac{e^{-\lambda} \lambda^{k}}{k!}, \quad k \in \set{0, 1, 2, \dots}
		\end{align*}
		\item Evaluamos explícitamente con $\lambda = 3$:
		\begin{align*}
			P(X = 0) &= e^{-3} \cdot \dfrac{3^{0}}{0!} = e^{-3} \approx 0.04979, \\
			P(X = 1) &= e^{-3} \cdot \dfrac{3^{1}}{1!} = 3 e^{-3} \approx 0.14936, \\
			P(X = 2) &= e^{-3} \cdot \dfrac{3^{2}}{2!} = \dfrac{9}{2} e^{-3} \approx 0.22404.
		\end{align*}
		Para la cola derecha:
		\begin{align*}
			P(X \ge 4) = 1 - P(X \le 3) = 1 - e^{-3}\left(1 + 3 + \dfrac{9}{2} + \dfrac{27}{6}\right) = 1 - 13 e^{-3} \approx 1 - 0.6472 = 0.3528.
		\end{align*}
	\end{enumerate}
\end{solucion}

\begin{solucion}
	\label{sol:3.7.2}
	Para el Problema \ref{prob:3.7.2}:
	\begin{enumerate}
		\item La PMF es $P(X = k) = e^{-2.5} \cdot 2.5^{k} / k!$. Calculamos las sumas sobre $k = 0, 1, \dots, 15$ (la cola más allá es despreciable):
		\begin{align*}
			\E[X] &= \sum_{k=0}^{15} k \cdot P(X=k) \approx 2.500, \\
			\E[X^{2}] &= \sum_{k=0}^{15} k^{2} \cdot P(X=k) \approx 8.750, \\
			\Var(X) &= \E[X^{2}] - (\E[X])^{2} \approx 8.750 - 6.250 = 2.500.
		\end{align*}
		\item La propiedad de equidispersión se confirma con tolerancia numérica: $\E[X] = \Var(X) = \lambda = 2.5$ (error $< 10^{-3}$). En control estadístico de procesos, la equidispersión implica que la carta de control $c$ (basada en Poisson) puede usar límites de control $3\sigma$ iguales a $\lambda \pm 3\sqrt{\lambda}$, simplificando el monitoreo industrial.
	\end{enumerate}
\end{solucion}

\begin{solucion}
	\label{sol:3.7.3}
	Para el Problema \ref{prob:3.7.3}:
	\begin{enumerate}
		\item Con $\lambda = 5$:
		\begin{align*}
			P(X \le 3) &= e^{-5}\left(\dfrac{5^{0}}{0!} + \dfrac{5^{1}}{1!} + \dfrac{5^{2}}{2!} + \dfrac{5^{3}}{3!}\right) = e^{-5}\left(1 + 5 + 12.5 + 20.833\overline{3}\right) \\
			&= e^{-5} \cdot 39.833\overline{3} \approx 0.006738 \cdot 39.833 \approx 0.2650.
		\end{align*}
		\item Para el complemento de la cola derecha:
		\begin{align*}
			P(X \ge 8) = 1 - P(X \le 7) = 1 - e^{-5} \sum_{k=0}^{7} \dfrac{5^{k}}{k!} \approx 1 - 0.867 = 0.1330.
		\end{align*}
		Por lo tanto, hay una probabilidad aproximada del $26.50\%$ de que lleguen a lo sumo $3$ pacientes y del $13.30\%$ de que lleguen al menos $8$ pacientes en un día típico.
	\end{enumerate}
\end{solucion}

\begin{solucion}
	\label{sol:3.7.4}
	Para el Problema \ref{prob:3.7.4}:
	\begin{enumerate}
		\item Bajo el modelo binomial exacto $Y \sim \text{Bin}(n = 2000, p = 0.001)$:
		\begin{align*}
			P(Y = 2) &= \comb{2000}{2}(0.001)^{2}(0.999)^{1998} \\
			&= \dfrac{2000 \cdot 1999}{2} \cdot 10^{-6} \cdot 0.999^{1998} \\
			&\approx 1\,999\,000 \cdot 10^{-6} \cdot 0.13534 \approx 0.27056.
		\end{align*}
		\item Bajo la aproximación de Poisson con $\lambda = np = 2$:
		\begin{align*}
			P(X = 2) = \dfrac{e^{-2} \cdot 2^{2}}{2!} = 2 e^{-2} \approx 2 \cdot 0.13534 \approx 0.27067.
		\end{align*}
		El error absoluto es $\abs{0.27056 - 0.27067} \approx 0.00011 < 5 \times 10^{-4}$.
		\item La aproximación de Poisson es excelente porque se cumplen las condiciones de la ley de eventos raros: $n = 2000 \gg 50$ y $p = 0.001 \ll 0.05$ con $np = 2 \le 5$. Bajo estos parámetros, los términos factoriales descendentes $(N-k)!/N!$ se simplifican, $(1-p)^{N-k} \to e^{-\lambda}$, y $\binom{N}{k} p^{k} \to \lambda^{k}/k!$.
	\end{enumerate}
\end{solucion}

\begin{solucion}
	\label{sol:3.7.5}
	Para el Problema \ref{prob:3.7.5}:
	\begin{enumerate}
		\item Con $\lambda = 8$:
		\begin{align*}
			P(X = 0) = e^{-8} \cdot \dfrac{8^{0}}{0!} = e^{-8} \approx 0.000335.
		\end{align*}
		\item La probabilidad de saturación se calcula mediante el complemento:
		\begin{align*}
			P(X \ge 15) = 1 - P(X \le 14) = 1 - \sum_{k=0}^{14} \dfrac{e^{-8} 8^{k}}{k!} \approx 1 - 0.9820 = 0.0180.
		\end{align*}
		\item Para el cuantil $q_{0.95}$, iteramos en `scipy.stats.poisson.cdf(q, mu=8)`:
		\begin{align*}
			P(X \le 11) &\approx 0.888, \quad P(X \le 12) \approx 0.936, \quad P(X \le 13) \approx 0.966.
		\end{align*}
		El menor $q$ tal que $P(X \le q) \ge 0.95$ es $q_{0.95} = 13$.
	\end{enumerate}
\end{solucion}

\begin{solucion}
	\label{sol:3.7.6}
	Para el Problema \ref{prob:3.7.6}:
	\begin{enumerate}
		\item Por la propiedad de escalamiento lineal del tiempo en un proceso de Poisson, el número de fallas en $4$ horas sigue una $\text{Pois}(\lambda' = 4 \cdot 0.5 = 2)$:
		\begin{align*}
			P(X \ge 1) = 1 - P(X = 0) = 1 - e^{-2} \approx 1 - 0.1353 = 0.8647.
		\end{align*}
		\item En el modelo binomial equivalente, se discretiza el turno en $n = 4$ ensayos horarios independientes con $p = 1 - e^{-0.5} \approx 0.3935$ (probabilidad de al menos una falla en una hora):
		\begin{align*}
			P(Y \ge 1) = 1 - (1-p)^{4} = 1 - (0.6065)^{4} \approx 1 - 0.1353 = 0.8647.
		\end{align*}
		Ambos enfoques producen exactamente el mismo resultado, validando la equivalencia operativa entre la formulación continua (Poisson) y la discretización Bernoulli horaria.
	\end{enumerate}
\end{solucion}

\begin{solucion}
	\label{sol:3.7.7}
	Para el Problema \ref{prob:3.7.7}:
	\begin{enumerate}
		\item Partimos de la identidad de la serie exponencial $e^{\lambda} = \sum_{k=0}^{\infty} \lambda^{k}/k!$ y la multiplicamos por $e^{-\lambda}$ para obtener la normalización de Poisson:
		\begin{align*}
			1 = e^{-\lambda} e^{\lambda} = \sum_{k=0}^{\infty} \dfrac{e^{-\lambda} \lambda^{k}}{k!} = \sum_{k=0}^{\infty} P(X = k).
		\end{align*}
		Derivamos término a término respecto a $\lambda$:
		\begin{align*}
			0 &= \dfrac{\partial}{\partial \lambda} \sum_{k=0}^{\infty} P(X = k) = \sum_{k=0}^{\infty} \dfrac{\partial}{\partial \lambda}\left(\dfrac{e^{-\lambda} \lambda^{k}}{k!}\right) = \sum_{k=0}^{\infty} \dfrac{e^{-\lambda} \lambda^{k-1}(\lambda - k)}{k!}.
		\end{align*}
		Multiplicamos por $\lambda$ y separamos el término $k = 0$:
		\begin{align*}
			0 = \sum_{k=0}^{\infty} \dfrac{e^{-\lambda} \lambda^{k}(k - \lambda)}{k!} = \sum_{k=0}^{\infty} k P(X = k) - \lambda \sum_{k=0}^{\infty} P(X = k) = \E[X] - \lambda.
		\end{align*}
		Por lo tanto $\E[X] = \lambda$.
		\item Para el segundo momento, derivamos dos veces y multiplicamos por $\lambda^{2}$:
		\begin{align*}
			\E[X(X-1)] = \sum_{k=0}^{\infty} k(k-1) P(X = k) = \lambda^{2}.
		\end{align*}
		Entonces:
		\begin{align*}
			\E[X^{2}] = \E[X(X-1)] + \E[X] = \lambda^{2} + \lambda.
		\end{align*}
		\item La varianza es:
		\begin{align*}
			\Var(X) = \E[X^{2}] - (\E[X])^{2} = (\lambda^{2} + \lambda) - \lambda^{2} = \lambda. \quad \blacksquare
		\end{align*}
	\end{enumerate}
\end{solucion}

\begin{solucion}
	\label{sol:3.7.8}
	Para el Problema \ref{prob:3.7.8}:
	\begin{enumerate}
		\item Por independencia, la convolución de las PMFs es:
		\begin{align*}
			P(Z = m) &= \sum_{k=0}^{m} P(X = k) P(Y = m-k) \\
			&= \sum_{k=0}^{m} \dfrac{e^{-\lambda_{1}} \lambda_{1}^{k}}{k!} \cdot \dfrac{e^{-\lambda_{2}} \lambda_{2}^{m-k}}{(m-k)!} \\
			&= e^{-(\lambda_{1} + \lambda_{2})} \sum_{k=0}^{m} \dfrac{\lambda_{1}^{k} \lambda_{2}^{m-k}}{k! (m-k)!} \\
			&= \dfrac{e^{-(\lambda_{1} + \lambda_{2})}}{m!} \sum_{k=0}^{m} \dfrac{m!}{k!(m-k)!} \lambda_{1}^{k} \lambda_{2}^{m-k}.
		\end{align*}
		Reconociendo el coeficiente binomial $\comb{m}{k} = m!/(k!(m-k)!)$ y aplicando el Teorema del Binomio de Newton:
		\begin{align*}
			P(Z = m) = \dfrac{e^{-(\lambda_{1} + \lambda_{2})}}{m!} (\lambda_{1} + \lambda_{2})^{m},
		\end{align*}
		que es exactamente la PMF de una $\text{Pois}(\lambda_{1} + \lambda_{2})$.
		\item La esperanza y varianza de $Z$ son:
		\begin{align*}
			\E[Z] = \E[X] + \E[Y] = \lambda_{1} + \lambda_{2}, \quad \Var(Z) = \Var(X) + \Var(Y) = \lambda_{1} + \lambda_{2}.
		\end{align*}
		Estos valores son consistentes con la PMF derivada, confirmando la coherencia algebraica del resultado. \qedhere
	\end{enumerate}
\end{solucion}

\begin{solucion}
	\label{sol:3.7.9}
	Para el Problema \ref{prob:3.7.9}:
	\begin{enumerate}
		\item Por definición de probabilidad condicional, independencia, y el resultado del Problema \ref{prob:3.7.8}:
		\begin{align*}
			P(X = k \mid X + Y = n) &= \dfrac{P(X = k) P(Y = n-k)}{P(X + Y = n)} \\
			&= \dfrac{\dfrac{e^{-\lambda_{1}} \lambda_{1}^{k}}{k!} \cdot \dfrac{e^{-\lambda_{2}} \lambda_{2}^{n-k}}{(n-k)!}}{\dfrac{e^{-(\lambda_{1}+\lambda_{2})} (\lambda_{1} + \lambda_{2})^{n}}{n!}} \\
			&= \dfrac{n!}{k! (n-k)!} \cdot \dfrac{\lambda_{1}^{k} \lambda_{2}^{n-k}}{(\lambda_{1} + \lambda_{2})^{n}} \\
			&= \comb{n}{k} \left(\dfrac{\lambda_{1}}{\lambda_{1} + \lambda_{2}}\right)^{k} \left(\dfrac{\lambda_{2}}{\lambda_{1} + \lambda_{2}}\right)^{n-k},
		\end{align*}
		que es la PMF de una $\text{Bin}(n, p)$ con $p = \lambda_{1}/(\lambda_{1} + \lambda_{2})$.
		\item Para $\lambda_{1} = 3$ y $\lambda_{2} = 5$, se tiene $p = 3/8 = 0.375$. Si $n = 8$ y $k = 3$:
		\begin{align*}
			P(X = 3 \mid X + Y = 8) &= \comb{8}{3} \left(\dfrac{3}{8}\right)^{3} \left(\dfrac{5}{8}\right)^{5} \\
			&= 56 \cdot \dfrac{27}{512} \cdot \dfrac{3125}{32768} \\
			&= 56 \cdot \dfrac{84\,375}{16\,777\,216} \approx 56 \cdot 0.00503 \approx 0.2815.
		\end{align*}
		Por lo tanto, dado que se atendieron en total $8$ solicitudes, la probabilidad de que el servidor $1$ haya atendido exactamente $3$ es aproximadamente $28.15\%$.
	\end{enumerate}
\end{solucion}

\begin{solucion}
	\label{sol:3.7.10}
	Para el Problema \ref{prob:3.7.10}:
	\begin{enumerate}
		\item En un proceso espacial de Poisson con intensidad $\lambda = 0.5$ defectos/$\text{m}^{2}$ y área total $A = 10\,\text{m}^{2}$, el número total de defectos sigue $X \sim \text{Pois}(\lambda A) = \text{Pois}(5)$. Evaluamos:
		\begin{align*}
			P(X = 0) &= e^{-5} \approx 0.00674, \\
			P(X \ge 8) &= 1 - P(X \le 7) \approx 1 - 0.8666 = 0.1334.
		\end{align*}
		\item Las regiones $A_{1} = 4\,\text{m}^{2}$ y $A_{2} = 6\,\text{m}^{2}$ son disjuntas, por lo que los conteos $X_{1}$ y $X_{2}$ son independientes con $X_{i} \sim \text{Pois}(\lambda A_{i})$:
		\begin{align*}
			X_{1} \sim \text{Pois}(0.5 \cdot 4) = \text{Pois}(2), \quad X_{2} \sim \text{Pois}(0.5 \cdot 6) = \text{Pois}(3).
		\end{align*}
		\item Por el resultado del Problema \ref{prob:3.7.9} aplicado con $\lambda_{1} = 2$ y $\lambda_{2} = 3$ (las medias efectivas de las sub-regiones), la distribución condicional es:
		\begin{align*}
			P(X_{1} = k \mid X_{1} + X_{2} = n) = \comb{n}{k} \left(\dfrac{2}{5}\right)^{k} \left(\dfrac{3}{5}\right)^{n-k} = \comb{n}{k} \left(\dfrac{A_{1}}{A}\right)^{k} \left(\dfrac{A_{2}}{A}\right)^{n-k}.
		\end{align*}
		Con $A_{1}/A = 4/10 = 0.4$, obtenemos $X_{1} \mid (X_{1} + X_{2} = n) \sim \text{Bin}(n, 0.4)$. Este resultado es intuitivo: dado un número fijo $n$ de defectos distribuidos al azar en un área total, la proporción que cae en una sub-región es proporcional a su área relativa. \qedhere
	\end{enumerate}
\end{solucion}

% ==============================================================================
% SECCIÓN 03.08: DISTRIBUCIÓN MULTINOMIAL Y ENSAYOS POLITÓMICOS
% ==============================================================================

\subsection*{Enunciados de los problemas --- Sección 03.08}

\subsubsection*{Nivel Fundamental}

\begin{problema}
	\label{prob:3.8.1}
	Considere un experimento aleatorio consistente en lanzar un dado cúbico justo $n = 6$ veces de manera independiente. Sea $\bm{X} = (X_{1}, X_{2}, \dots, X_{6})$ el vector aleatorio que cuenta el número de veces que aparece cada cara, donde $X_{i}$ es el número de veces que sale la cara $i$ en los $6$ lanzamientos.
	\begin{enumerate}
		\item Especifique formalmente la distribución del vector $\bm{X}$ utilizando la distribución multinomial con parámetros $n = 6$ y probabilidades $p_{i} = 1/6$ para todo $i$.
		\item Calcule explícitamente la probabilidad de que cada una de las 6 caras aparezca exactamente una vez.
	\end{enumerate}
\end{problema}

\begin{problema}
	\label{prob:3.8.2}
	Una encuesta de opinión política se aplica a $n = 100$ votantes seleccionados al azar, donde cada votante elige uno de tres candidatos con probabilidades $p_{1} = 0.45$, $p_{2} = 0.35$ y $p_{3} = 0.20$. Sea $\bm{X} = (X_{1}, X_{2}, X_{3})$ el número de votos获得的 por cada candidato.
	\begin{enumerate}
		\item Calcule la probabilidad de que el candidato 1 reciba exactamente $X_{1} = 50$ votos, el candidato 2 reciba $X_{2} = 30$ y el candidato 3 reciba $X_{3} = 20$.
		\item Calcule la probabilidad de que el candidato 1 obtenga mayoría absoluta ($X_{1} \ge 51$).
	\end{enumerate}
\end{problema}

\begin{problema}
	\label{prob:3.8.3}
	Una empresa de manufactura produce lotes de chips en los que cada chip se clasifica en tres categorías mutuamente excluyentes: ``Premium'' ($p_{1} = 0.70$), ``Estándar'' ($p_{2} = 0.25$) y ``Defectuoso'' ($p_{3} = 0.05$). En un lote de $n = 10$ chips, sea $\bm{X} = (X_{1}, X_{2}, X_{3})$ el vector de conteos por categoría.
	\begin{enumerate}
		\item Calcule el número total de configuraciones distintas en las que pueden distribuirse los 10 chips en las 3 categorías: $\binom{10 + 3 - 1}{3 - 1} = \binom{12}{2}$.
		\item Calcule la probabilidad de que en un lote de 10 chips se obtengan exactamente 7 Premium, 2 Estándar y 1 Defectuoso.
	\end{enumerate}
\end{problema}

\subsubsection*{Nivel Operativo}

\begin{problema}
	\label{prob:3.8.4}
	Considere un experimento con $n = 100$ ensayos independientes donde cada ensayo puede resultar en una de tres categorías con probabilidades $p_{1} = 0.5$, $p_{2} = 0.3$ y $p_{3} = 0.2$.
	\begin{enumerate}
		\item Demuestre que la marginal $X_{1} \sim \text{Bin}(100, 0.5)$ sumando sobre todos los valores posibles de $X_{2}$ y $X_{3}$.
		\item Calcule $\E[X_{1}]$, $\Var(X_{1})$ y la probabilidad $P(X_{1} = 50)$.
	\end{enumerate}
\end{problema}

\begin{problema}
	\label{prob:3.8.5}
	En un estudio de genómica de poblaciones, se analizan $n = 200$ sitios genéticos independientes donde cada sitio puede presentar uno de tres genotipos con probabilidades $p_{1} = 0.36$ (homocigoto dominante), $p_{2} = 0.48$ (heterocigoto) y $p_{3} = 0.16$ (homocigoto recesivo) bajo equilibrio de Hardy-Weinberg.
	\begin{enumerate}
		\item Calcule la esperanza y la varianza del número de sitios homocigotos recesivos $X_{3}$.
		\item Calcule la probabilidad de observar al menos 40 sitios homocigotos recesivos, $P(X_{3} \ge 40)$.
	\end{enumerate}
\end{problema}

\begin{problema}
	\label{prob:3.8.6}
	Una compañía de seguros clasifica su cartera de $n = 1000$ pólizas en tres categorías de riesgo: bajo ($p_{1} = 0.60$), medio ($p_{2} = 0.30$) y alto ($p_{3} = 0.10$). Sea $\bm{X} = (X_{1}, X_{2}, X_{3})$ el vector de conteos.
	\begin{enumerate}
		\item Simule computacionalmente (usando `numpy.random.multinomial`) una realización de $\bm{X}$ y reporte los conteos obtenidos.
		\item Repita la simulación 250{,}000 veces y calcule las medias y varianzas empíricas de $X_{1}, X_{2}, X_{3}$ para verificar que coincidan con los valores teóricos.
	\end{enumerate}
\end{problema}

\subsubsection*{Nivel Analítico}

\begin{problema}
	\label{prob:3.8.7}
	Sea $\bm{X} = (X_{1}, X_{2}, \dots, X_{k}) \sim \text{Mult}(n, p_{1}, p_{2}, \dots, p_{k})$ un vector aleatorio con distribución multinomial.
	\begin{enumerate}
		\item Demuestre analíticamente que la covarianza entre dos componentes distintas es:
		\begin{align*}
			\cov(X_{i}, X_{j}) = -n p_{i} p_{j}, \quad \text{para } i \neq j.
		\end{align*}
		\item Interprete el signo negativo de la covarianza: ¿por qué un exceso en la categoría $i$ implica necesariamente un déficit en la categoría $j$?
	\end{enumerate}
\end{problema}

\begin{problema}
	\label{prob:3.8.8}
	Sea $\bm{X} = (X_{1}, X_{2}, X_{3}) \sim \text{Mult}(n, p_{1}, p_{2}, p_{3})$ con $n$ ensayos independientes.
	\begin{enumerate}
		\item Deduzca analíticamente el momento factorial conjunto $\E[X_{1}(X_{1}-1)]$ aplicando propiedades de la expansión del multinomio.
		\item Generalice el resultado al caso de $k$ categorías: demuestre que $\E[X_{i}(X_{i}-1)] = n(n-1)p_{i}^{2}$ y $\E[X_{i} X_{j}] = n(n-1)p_{i} p_{j}$ para $i \neq j$.
	\end{enumerate}
\end{problema}

\subsubsection*{Nivel Desafiante}

\begin{problema}
	\label{prob:3.8.9}
	Sea $\bm{X} = (X_{1}, X_{2}, X_{3}, X_{4}) \sim \text{Mult}(n, p_{1}, p_{2}, p_{3}, p_{4})$ con cuatro categorías mutuamente excluyentes.
	\begin{enumerate}
		\item Demuestre que, condicionada al total $X_{3} + X_{4} = m$, la distribución conjunta del sub-vector $(X_{3}, X_{4})$ colapsa a una distribución binomial $\text{Bin}(m, p_{3}/(p_{3} + p_{4}))$ en el conteo de $X_{3}$.
		\item Aplique el resultado a un caso con $n = 100$, $p_{1} = 0.4$, $p_{2} = 0.3$, $p_{3} = 0.2$, $p_{4} = 0.1$: si en una muestra de 100 ensayos se observa que $X_{3} + X_{4} = 30$, calcule la probabilidad de que $X_{3} = 18$ y $X_{4} = 12$.
	\end{enumerate}
\end{problema}

\begin{problema}
	\label{prob:3.8.10}
	En un estudio de mercado se quiere evaluar la independencia entre la preferencia de marca (A, B, C) y el grupo etario (joven, adulto) en una muestra de $n = 500$ consumidores. Sea $X_{ij}$ el conteo en la celda $(i, j)$ de la tabla de contingencia $3 \times 2$.
	\begin{enumerate}
		\item Plantee la hipótesis nula $H_{0}$ de independencia entre marca y grupo etario, especificando las probabilidades conjuntas $p_{ij} = p_{i}^{\text{marca}} \cdot p_{j}^{\text{edad}}$ bajo $H_{0}$.
		\item Bajo $H_{0}$ con probabilidades marginales $p_{1}^{\text{marca}} = 0.4$, $p_{2}^{\text{marca}} = 0.35$, $p_{3}^{\text{marca}} = 0.25$ y $p_{1}^{\text{edad}} = 0.6$ (joven), $p_{2}^{\text{edad}} = 0.4$ (adulto), calcule el valor esperado de cada celda: $E_{ij} = n p_{i}^{\text{marca}} p_{j}^{\text{edad}}$.
		\item Demuestre que, condicionada a los marginales fijos de fila y columna, la distribución de las celdas sigue una distribución multinomial condicional análoga al Problema \ref{prob:3.8.9}.
	\end{enumerate}
\end{problema}

\subsection*{Sugerencias de los problemas --- Sección 03.08}

\begin{sugerencia}
	\label{sug:3.8.1}
	Para el Problema \ref{prob:3.8.1}, use la PMF multinomial $P(\bm{X} = \bm{x}) = \dfrac{n!}{x_{1}!\cdots x_{k}!} p_{1}^{x_{1}} \cdots p_{k}^{x_{k}}$ con $n=6$, $k=6$, $p_{i} = 1/6$ y $x_{i} = 1$ para todo $i$. El coeficiente multinomial vale $6!/(1!)^{6} = 720$.
\end{sugerencia}

\begin{sugerencia}
	\label{sug:3.8.2}
	Para el Problema \ref{prob:3.8.2}, evalúe el coeficiente multinomial $\binom{100}{50, 30, 20} = \dfrac{100!}{50! \cdot 30! \cdot 20!}$ con `scipy.special.multinomial`. Para la mayoría absoluta, sume $P(X_{1} = k)$ sobre $k \ge 51$ marginalizando sobre las otras categorías.
\end{sugerencia}

\begin{sugerencia}
	\label{sug:3.8.3}
	Para el Problema \ref{prob:3.8.3}, use que el número de formas de distribuir $n$ objetos en $k$ categorías con orden interno es el coeficiente multinomial $\dfrac{n!}{x_{1}! \cdots x_{k}!}$, no el de estrellas y barras. El coeficiente multinomial de $n=10$ en $(7, 2, 1)$ vale $10!/(7! \cdot 2! \cdot 1!) = 360$.
\end{sugerencia}

\begin{sugerencia}
	\label{sug:3.8.4}
	Para el Problema \ref{prob:3.8.4}, la marginalización se realiza sumando $P(X_{1} = x_{1}, X_{2} = x_{2}, X_{3} = x_{3})$ sobre todos los $x_{2} \in \set{0, \dots, n - x_{1}}$ con $x_{3} = n - x_{1} - x_{2}$. Use el Teorema del Binomio: $\sum_{x_{2}=0}^{n-x_{1}} \binom{n - x_{1}}{x_{2}} (p_{2}/(1-p_{1}))^{x_{2}} = (1-p_{1})^{n-x_{1}} \cdot (1/(1-p_{1}))^{n-x_{1}} = 1$.
\end{sugerencia}

\begin{sugerencia}
	\label{sug:3.8.5}
	Para el Problema \ref{prob:3.8.5}, use que $X_{3}$ es marginalmente binomial: $X_{3} \sim \text{Bin}(n = 200, p_{3} = 0.16)$. Entonces $\E[X_{3}] = np_{3} = 32$, $\Var(X_{3}) = np_{3}(1-p_{3}) = 26.88$, y $P(X_{3} \ge 40) = 1 - \text{stats.binom.cdf}(39, 200, 0.16)$.
\end{sugerencia}

\begin{sugerencia}
	\label{sug:3.8.6}
	Para el Problema \ref{prob:3.8.6}, use `np.random.seed(42)` y `np.random.multinomial(n=1000, pvals=[0.6, 0.3, 0.1])` para una sola realización. Para la simulación masiva, genere un array de $250{,}000$ realizaciones y calcule `np.mean(axis=0)` y `np.var(axis=0, ddof=1)` para obtener las medias y varianzas empíricas.
\end{sugerencia}

\begin{sugerencia}
	\label{sug:3.8.7}
	Para el Problema \ref{prob:3.8.7}, use la representación $X_{i} = \sum_{r=1}^{n} I_{r}^{(i)}$ donde $I_{r}^{(i)} = 1$ si el ensayo $r$ produce categoría $i$. Entonces $\cov(X_{i}, X_{j}) = \sum_{r=1}^{n} \sum_{s=1}^{n} \cov(I_{r}^{(i)}, I_{s}^{(j)})$. Para $r = s$, $\cov(I_{r}^{(i)}, I_{r}^{(j)}) = -p_{i} p_{j}$ por exclusión mutua; para $r \neq s$, la covarianza es cero por independencia entre ensayos. Por lo tanto $\cov(X_{i}, X_{j}) = -n p_{i} p_{j}$.
\end{sugerencia}

\begin{sugerencia}
	\label{sug:3.8.8}
	Para el Problema \ref{prob:3.8.8}, exprese la suma multinomial $\sum_{x_{1}+x_{2}+x_{3}=n} \dfrac{n!}{x_{1}! x_{2}! x_{3}!} p_{1}^{x_{1}} p_{2}^{x_{2}} p_{3}^{x_{3}}$. Para $\E[X_{1}(X_{1}-1)]$, multiplique el sumando por $x_{1}(x_{1}-1)$ y use la identidad $x_{1}(x_{1}-1) \dfrac{n!}{x_{1}!} = n(n-1) \dfrac{(n-2)!}{(x_{1}-2)!}$, reconociendo una nueva distribución multinomial con $n-2$ ensayos.
\end{sugerencia}

\begin{sugerencia}
	\label{sug:3.8.9}
	Para el Problema \ref{prob:3.8.9}, use probabilidad condicional y la PMF multinomial: $P(X_{3} = k \mid X_{3} + X_{4} = m) = \dfrac{P(X_{3} = k, X_{4} = m-k)}{P(X_{3} + X_{4} = m)}$. El numerador es el término de la multinomial con $(n, p_{1}, p_{2}, p_{3}, p_{4})$ y $X_{1}, X_{2}$ absorbidos por marginalización. El denominador es el término binomial con $n$ ensayos y probabilidad $p_{3} + p_{4}$.
\end{sugerencia}

\begin{sugerencia}
	\label{sug:3.8.10}
	Para el Problema \ref{prob:3.8.10}, bajo $H_{0}$, las celdas $\{X_{ij}\}$ siguen una distribución multinomial con $n = 500$ ensayos y probabilidades $p_{ij} = p_{i}^{\text{marca}} p_{j}^{\text{edad}}$. Los valores esperados son $E_{ij} = 500 p_{ij}$. La distribución condicional de una celda dado los marginales fijos es hipergeométrica multivariada (generalización de la prueba exacta de Fisher).
\end{sugerencia}

\subsection*{Soluciones de los problemas --- Sección 03.08}

\begin{solucion}
	\label{sol:3.8.1}
	Para el Problema \ref{prob:3.8.1}:
	\begin{enumerate}
		\item El vector $\bm{X} = (X_{1}, \dots, X_{6})$ sigue una distribución multinomial con parámetros $n = 6$, $k = 6$ categorías y probabilidades $p_{i} = 1/6$ para todo $i$:
		\begin{align*}
			P(X_{1} = x_{1}, \dots, X_{6} = x_{6}) = \dfrac{6!}{x_{1}! \cdots x_{6}!} \left(\dfrac{1}{6}\right)^{6}, \quad \sum_{i=1}^{6} x_{i} = 6.
		\end{align*}
		\item Para el evento de que cada cara aparezca exactamente una vez, tenemos $x_{i} = 1$ para todo $i$:
		\begin{align*}
			P(X_{1} = 1, \dots, X_{6} = 1) &= \dfrac{6!}{1! \cdot 1! \cdot 1! \cdot 1! \cdot 1! \cdot 1!} \cdot \left(\dfrac{1}{6}\right)^{6} \\
			&= 720 \cdot \dfrac{1}{46{,}656} = \dfrac{720}{46{,}656} \approx 0.01543.
		\end{align*}
		Por lo tanto, la probabilidad de obtener las 6 caras diferentes en 6 lanzamientos es aproximadamente $1.54\%$.
	\end{enumerate}
\end{solucion}

\begin{solucion}
	\label{sol:3.8.2}
	Para el Problema \ref{prob:3.8.2}:
	\begin{enumerate}
		\item Con $\bm{X} \sim \text{Mult}(n = 100, p = (0.45, 0.35, 0.20))$ y vector de conteos $(50, 30, 20)$:
		\begin{align*}
			P(X_{1} = 50, X_{2} = 30, X_{3} = 20) &= \binom{100}{50, 30, 20} \cdot (0.45)^{50} \cdot (0.35)^{30} \cdot (0.20)^{20} \\
			&\approx 2.34 \times 10^{27} \cdot 9.13 \times 10^{-18} \cdot 1.93 \times 10^{-14} \cdot 1.05 \times 10^{-14} \\
			&\approx 4.32 \times 10^{-18}.
		\end{align*}
		\item Para la mayoría absoluta del candidato 1, marginalizamos sobre $X_{2}$ y $X_{3}$ con la restricción $X_{2} + X_{3} = 100 - X_{1}$:
		\begin{align*}
			P(X_{1} \ge 51) &= \sum_{x_{1}=51}^{100} \binom{100}{x_{1}} (0.45)^{x_{1}} (0.55)^{100-x_{1}} \\
			&\approx 1 - \Phi\left(\dfrac{50.5 - 45}{\sqrt{100 \cdot 0.45 \cdot 0.55}}\right) \\
			&\approx 1 - \Phi(1.099) \approx 1 - 0.864 = 0.136.
		\end{align*}
		La probabilidad de mayoría absoluta del candidato 1 es aproximadamente $13.6\%$.
	\end{enumerate}
\end{solucion}

\begin{solucion}
	\label{sol:3.8.3}
	Para el Problema \ref{prob:3.8.3}:
	\begin{enumerate}
		\item El número de configuraciones distintas en que se pueden distribuir $n = 10$ chips en $k = 3$ categorías ordenadas (cada chip pertenece a una categoría específica) está dado por el coeficiente multinomial cuando se enumeran los chips:
		\begin{align*}
			\sum_{x_{1}+x_{2}+x_{3}=10} \binom{10}{x_{1}, x_{2}, x_{3}} = 3^{10} = 59{,}049.
		\end{align*}
		Esto se debe a que cada uno de los 10 chips puede asignarse a 3 categorías independientemente. El número de estrellas y barras $\binom{12}{2} = 66$ cuenta únicamente las particiones sin identidad de chips, no las configuraciones etiquetadas.
		\item Para la probabilidad específica del conteo $(7, 2, 1)$:
		\begin{align*}
			P(X_{1} = 7, X_{2} = 2, X_{3} = 1) &= \binom{10}{7, 2, 1} (0.70)^{7} (0.25)^{2} (0.05)^{1} \\
			&= \dfrac{10!}{7! \cdot 2! \cdot 1!} \cdot (0.0824) \cdot (0.0625) \cdot (0.05) \\
			&= 360 \cdot 2.575 \times 10^{-4} \approx 0.0927.
		\end{align*}
		La probabilidad es aproximadamente $9.27\%$.
	\end{enumerate}
\end{solucion}

\begin{solucion}
	\label{sol:3.8.4}
	Para el Problema \ref{prob:3.8.4}:
	\begin{enumerate}
		\item Marginalizamos la distribución multinomial sobre $X_{2}$ y $X_{3}$ con la restricción $X_{2} + X_{3} = n - X_{1}$:
		\begin{align*}
			P(X_{1} = x_{1}) &= \sum_{x_{2}=0}^{n-x_{1}} P(X_{1} = x_{1}, X_{2} = x_{2}, X_{3} = n - x_{1} - x_{2}) \\
			&= \sum_{x_{2}=0}^{n-x_{1}} \dfrac{n!}{x_{1}! x_{2}! (n - x_{1} - x_{2})!} p_{1}^{x_{1}} p_{2}^{x_{2}} p_{3}^{n-x_{1}-x_{2}} \\
			&= \dfrac{n!}{x_{1}! (n-x_{1})!} p_{1}^{x_{1}} (p_{2} + p_{3})^{n-x_{1}} \sum_{x_{2}=0}^{n-x_{1}} \binom{n-x_{1}}{x_{2}} \left(\dfrac{p_{2}}{p_{2}+p_{3}}\right)^{x_{2}} \left(\dfrac{p_{3}}{p_{2}+p_{3}}\right)^{n-x_{1}-x_{2}}.
		\end{align*}
		La suma final es exactamente $1$ por el Teorema del Binomio, por lo que:
		\begin{align*}
			P(X_{1} = x_{1}) = \binom{n}{x_{1}} p_{1}^{x_{1}} (1 - p_{1})^{n - x_{1}} \implies X_{1} \sim \text{Bin}(n, p_{1}).
		\end{align*}
		\item Con $n = 100$ y $p_{1} = 0.5$:
		\begin{align*}
			\E[X_{1}] &= 100 \cdot 0.5 = 50, \quad \Var(X_{1}) = 100 \cdot 0.5 \cdot 0.5 = 25, \\
			P(X_{1} = 50) &= \binom{100}{50} (0.5)^{100} \approx \dfrac{1}{\sqrt{2\pi \cdot 25}} \approx 0.0798.
		\end{align*}
	\end{enumerate}
\end{solucion}

\begin{solucion}
	\label{sol:3.8.5}
	Para el Problema \ref{prob:3.8.5}:
	\begin{enumerate}
		\item Por el resultado del Problema \ref{prob:3.8.4}, $X_{3} \sim \text{Bin}(n = 200, p_{3} = 0.16)$ es marginalmente binomial:
		\begin{align*}
			\E[X_{3}] &= np_{3} = 200 \cdot 0.16 = 32, \\
			\Var(X_{3}) &= np_{3}(1 - p_{3}) = 200 \cdot 0.16 \cdot 0.84 = 26.88.
		\end{align*}
		\item Para la cola derecha:
		\begin{align*}
			P(X_{3} \ge 40) &= 1 - P(X_{3} \le 39) = 1 - \text{stats.binom.cdf}(39, 200, 0.16) \\
			&\approx 1 - 0.910 = 0.090.
		\end{align*}
		Por lo tanto, la probabilidad de observar al menos 40 sitios homocigotos recesivos es aproximadamente $9.0\%$.
	\end{enumerate}
\end{solucion}

\begin{solucion}
	\label{sol:3.8.6}
	Para el Problema \ref{prob:3.8.6}:
	\begin{enumerate}
		\item Con `np.random.seed(42)` y `np.random.multinomial(1000, [0.6, 0.3, 0.1])`, una realización típica es, por ejemplo, $\bm{X} = (598, 305, 97)$.
		\item Con 250{,}000 simulaciones, los momentos empíricos convergen a los teóricos:
		\begin{align*}
			\bar{X}_{1} &\approx 600.0, \quad s_{X_{1}}^{2} \approx 240.0 \quad (\text{teórico: } \E = 600, \Var = 240) \\
			\bar{X}_{2} &\approx 300.0, \quad s_{X_{2}}^{2} \approx 210.0 \quad (\text{teórico: } \E = 300, \Var = 210) \\
			\bar{X}_{3} &\approx 100.0, \quad s_{X_{3}}^{2} \approx 90.0 \quad (\text{teórico: } \E = 100, \Var = 90)
		\end{align*}
		donde $\Var(X_{i}) = np_{i}(1-p_{i})$. Adicionalmente, las covarianzas empíricas deberían coincidir con $\cov(X_{i}, X_{j}) = -np_{i}p_{j}$: por ejemplo, $\widehat{\cov}(X_{1}, X_{2}) \approx -1000 \cdot 0.6 \cdot 0.3 = -180$.
	\end{enumerate}
\end{solucion}

\begin{solucion}
	\label{sol:3.8.7}
	Para el Problema \ref{prob:3.8.7}:
	\begin{enumerate}
		\item Representamos cada componente como suma de indicadores: $X_{i} = \sum_{r=1}^{n} I_{r}^{(i)}$ donde $I_{r}^{(i)} = 1$ si el ensayo $r$ resulta en categoría $i$. Por bilinealidad de la covarianza:
		\begin{align*}
			\cov(X_{i}, X_{j}) &= \sum_{r=1}^{n} \sum_{s=1}^{n} \cov(I_{r}^{(i)}, I_{s}^{(j)}).
		\end{align*}
		Si $r \neq s$, los ensayos son independientes y $\cov(I_{r}^{(i)}, I_{s}^{(j)}) = 0$. Si $r = s$, la covarianza es:
		\begin{align*}
			\cov(I_{r}^{(i)}, I_{r}^{(j)}) &= \E[I_{r}^{(i)} I_{r}^{(j)}] - \E[I_{r}^{(i)}] \E[I_{r}^{(j)}].
		\end{align*}
		Por la exclusión mutua, no es posible que un ensayo produzca simultáneamente las categorías $i$ y $j$ con $i \neq j$, así que $I_{r}^{(i)} I_{r}^{(j)} = 0$ y $\E[I_{r}^{(i)} I_{r}^{(j)}] = 0$. Entonces:
		\begin{align*}
			\cov(I_{r}^{(i)}, I_{r}^{(j)}) = 0 - p_{i} p_{j} = -p_{i} p_{j}.
		\end{align*}
		Sumando sobre los $n$ términos con $r = s$:
		\begin{align*}
			\cov(X_{i}, X_{j}) = n \cdot (-p_{i} p_{j}) = -n p_{i} p_{j} \quad \text{para } i \neq j. \quad \blacksquare
		\end{align*}
		\item El signo negativo se interpreta por la restricción $\sum_{i} X_{i} = n$: si el conteo de la categoría $i$ excede su valor esperado $np_{i}$ (es decir, hay ``exceso'' en $i$), entonces necesariamente algún conteo $X_{j}$ con $j \neq i$ debe estar por debajo de su valor esperado, ya que el total permanece fijo en $n$. Esta es la firma algebraica de la restricción lineal en el soporte del vector multinomial.
	\end{enumerate}
\end{solucion}

\begin{solucion}
	\label{sol:3.8.8}
	Para el Problema \ref{prob:3.8.8}:
	\begin{enumerate}
		\item Para $\E[X_{1}(X_{1}-1)]$, multiplicamos el sumando multinomial por $x_{1}(x_{1}-1)$:
		\begin{align*}
			\E[X_{1}(X_{1}-1)] &= \sum_{x_{1}+x_{2}+x_{3}=n} x_{1}(x_{1}-1) \cdot \dfrac{n!}{x_{1}! x_{2}! x_{3}!} p_{1}^{x_{1}} p_{2}^{x_{2}} p_{3}^{x_{3}}.
		\end{align*}
		Usando la identidad $x_{1}(x_{1}-1) \cdot \dfrac{n!}{x_{1}!} = n(n-1) \cdot \dfrac{(n-2)!}{(x_{1}-2)!}$, factorizamos:
		\begin{align*}
			\E[X_{1}(X_{1}-1)] &= n(n-1) p_{1}^{2} \sum_{x_{1}+x_{2}+x_{3}=n} \dfrac{(n-2)!}{(x_{1}-2)! x_{2}! x_{3}!} p_{1}^{x_{1}-2} p_{2}^{x_{2}} p_{3}^{x_{3}}.
		\end{align*}
		Con el cambio de variable $y_{1} = x_{1} - 2$ y reconociendo la suma como una distribución multinomial con $n - 2$ ensayos:
		\begin{align*}
			\E[X_{1}(X_{1}-1)] = n(n-1) p_{1}^{2}.
		\end{align*}
		\item Para $k$ categorías, análogamente:
		\begin{align*}
			\E[X_{i}(X_{i}-1)] &= n(n-1) p_{i}^{2}, \\
			\E[X_{i} X_{j}] &= n(n-1) p_{i} p_{j} \quad (i \neq j).
		\end{align*}
		Combinando ambos resultados: $\Var(X_{i}) = \E[X_{i}(X_{i}-1)] + \E[X_{i}] - (\E[X_{i}])^{2} = n(n-1)p_{i}^{2} + np_{i} - n^{2}p_{i}^{2} = np_{i}(1-p_{i})$, confirmando la varianza binomial marginal. \qedhere
	\end{enumerate}
\end{solucion}

\begin{solucion}
	\label{sol:3.8.9}
	Para el Problema \ref{prob:3.8.9}:
	\begin{enumerate}
		\item Por probabilidad condicional, marginalización y la independencia de la suma $X_{3} + X_{4}$ respecto al complemento:
		\begin{align*}
			P(X_{3} = k \mid X_{3} + X_{4} = m) &= \dfrac{P(X_{3} = k, X_{4} = m-k)}{P(X_{3} + X_{4} = m)} \\
			&= \dfrac{\binom{n}{k, m-k, n-m} p_{3}^{k} p_{4}^{m-k} (1 - p_{3} - p_{4})^{n-m}}{\binom{n}{m} (p_{3} + p_{4})^{m} (1 - p_{3} - p_{4})^{n-m}} \\
			&= \binom{m}{k} \left(\dfrac{p_{3}}{p_{3} + p_{4}}\right)^{k} \left(\dfrac{p_{4}}{p_{3} + p_{4}}\right)^{m-k}.
		\end{align*}
		Esta es la PMF de una distribución $\text{Bin}(m, p_{3}/(p_{3}+p_{4}))$.
		\item Con $n = 100$, $p_{3} = 0.2$, $p_{4} = 0.1$, $m = 30$, $k = 18$:
		\begin{align*}
			P(X_{3} = 18 \mid X_{3} + X_{4} = 30) &= \binom{30}{18} \left(\dfrac{0.2}{0.3}\right)^{18} \left(\dfrac{0.1}{0.3}\right)^{12} \\
			&= 86\,493\,225 \cdot (2/3)^{18} \cdot (1/3)^{12} \\
			&\approx 86\,493\,225 \cdot 6.76 \times 10^{-4} \cdot 1.87 \times 10^{-6} \\
			&\approx 0.1094.
		\end{align*}
		Por lo tanto, dado que $X_{3} + X_{4} = 30$, la probabilidad de que $X_{3} = 18$ y $X_{4} = 12$ es aproximadamente $10.94\%$.
	\end{enumerate}
\end{solucion}

\begin{solucion}
	\label{sol:3.8.10}
	Para el Problema \ref{prob:3.8.10}:
	\begin{enumerate}
		\item Bajo $H_{0}$ de independencia, las probabilidades conjuntas factorizan: $p_{ij} = p_{i}^{\text{marca}} \cdot p_{j}^{\text{edad}}$. Las celdas siguen una distribución multinomial con $n = 500$ ensayos y $6$ categorías (3 marcas $\times$ 2 grupos etarios). La hipótesis nula se evalúa comparando conteos observados $X_{ij}$ con valores esperados $E_{ij} = 500 p_{ij}$.
		\item Calculamos los valores esperados de las 6 celdas con las marginales dadas:
		\begin{align*}
			E_{11} = 500 \cdot 0.4 \cdot 0.6 = 120, \quad E_{12} &= 500 \cdot 0.4 \cdot 0.4 = 80, \\
			E_{21} = 500 \cdot 0.35 \cdot 0.6 = 105, \quad E_{22} &= 500 \cdot 0.35 \cdot 0.4 = 70, \\
			E_{31} = 500 \cdot 0.25 \cdot 0.6 = 75, \quad E_{32} &= 500 \cdot 0.25 \cdot 0.4 = 50.
		\end{align*}
		\item Condicionada a los marginales de fila $\bm{X}_{i \bullet} = (X_{i1} + X_{i2})$ y de columna $\bm{X}_{\bullet j} = (X_{1j} + X_{2j} + X_{3j})$, la distribución conjunta de las celdas es la \emph{multinomial condicional} (o distribución hipergeométrica multivariada): cada celda $X_{ij}$ dado los marginales sigue una distribución hipergeométrica con parámetros poblacionales iguales a los marginales observados. Este es el análogo multivariado de la prueba exacta de Fisher para tablas $r \times c$. \qedhere
	\end{enumerate}
\end{solucion}
